\section{Proofs for Section~\ref{sec:algebraic}}
\label{app:algebraic}

This appendix proves every result of Section~\ref{sec:algebraic}. Norms of
vectors are Euclidean and norms of matrices are operator norms unless the
Frobenius norm $\norm\cdot_F$ is indicated. For a rational number we use its
reduced fraction. A \emph{quadratic} is a polynomial of degree at most two.
Imported results are stated with their exact contracts where they are used;
the algebraic-geometric ones are collected in
Appendix~\ref{app:algebraic-imports}.

\subsection{Canonical Hessian Grams}
\label{app:algebraic-gram}

\begin{proof}[Proof of Lemma~\ref{lem:algebraic-gram}]
\emph{(a)} First let $f=g^2$ with $g(x)=g_0+b^{\mathsf T}x+x^{\mathsf T}Tx$.
Then $\nabla g(x)=b+2Tx$ and $\nabla^2g=2T$, so
\[
 v^{\mathsf T}\nabla^2(g^2)(x)v
 =2\bigl(v^{\mathsf T}\nabla g(x)\bigr)^2+2g(x)\,v^{\mathsf T}\nabla^2g\,v
 =2\bigl(b^{\mathsf T}v+2x^{\mathsf T}Tv\bigr)^2
  +4\bigl(g_0+b^{\mathsf T}x+x^{\mathsf T}Tx\bigr)v^{\mathsf T}Tv .
\]
On the other hand, with $w=w(x,v)$, the three blocks of $\Gamma(g)$
contribute
\[
\begin{aligned}
 v^{\mathsf T}(2bb^{\mathsf T}+4g_0T)v
   &=2(b^{\mathsf T}v)^2+4g_0\,v^{\mathsf T}Tv,\\
 2v^{\mathsf T}\Delta(b,T)(x\otimes v)
   &=2\sum_{i,k,j}v_i\bigl(2b_kT_{ij}+4b_iT_{kj}\bigr)x_kv_j\\
   &=4(b^{\mathsf T}x)(v^{\mathsf T}Tv)+8(b^{\mathsf T}v)(x^{\mathsf T}Tv),\\
 (x\otimes v)^{\mathsf T}\Xi(T)(x\otimes v)
   &=8\Bigl(\sum_{k,j}x_kv_jT_{kj}\Bigr)^2
     +4\sum_{k,j,k',j'}x_kv_jT_{kk'}T_{jj'}x_{k'}v_{j'}\\
   &=8(x^{\mathsf T}Tv)^2+4(x^{\mathsf T}Tx)(v^{\mathsf T}Tv).
\end{aligned}
\]
Their sum equals the previous display. Thus
$v^{\mathsf T}\nabla^2(g^2)(x)v=w^{\mathsf T}\Gamma(g)w$ for every quadratic
$g$. Both sides are quadratic functions of the coefficient vector of $g$.
Polarizing, $v^{\mathsf T}\nabla^2(g_kg_l)(x)v=w^{\mathsf
T}\widehat\Gamma(g_k,g_l)w$, because
$2g_kg_l=(g_k+g_l)^2-g_k^2-g_l^2$. Summing with the weights $S_{kl}$ proves
(a).

\emph{(b)} Again it suffices to treat one quadratic: both sides of the
claimed identity are quadratic in the coefficients of $g$, the recentring
$g\mapsto g(\cdot+c)$ is linear in those coefficients, and the polarized and
weighted identities follow as in (a). Put $A_c=c\otimes I_n$, so that
$(A_cv)_{(k,j)}=c_kv_j$ and $(A_c^{\mathsf T}y)_i=\sum_kc_ky_{(k,i)}$. For
every symmetric $T$ and every vector $b$, direct contraction of indices gives
\begin{align}
 A_c^{\mathsf T}\Xi(T)&=\Delta(2Tc,T),
 \label{eq:algebraic-contract-1}\\
 A_c^{\mathsf T}\Xi(T)A_c&=8(Tc)(Tc)^{\mathsf T}+4(c^{\mathsf T}Tc)\,T,
 \label{eq:algebraic-contract-2}\\
 \Delta(b,T)A_c+A_c^{\mathsf T}\Delta(b,T)^{\mathsf T}
  &=4(b^{\mathsf T}c)\,T+4b(Tc)^{\mathsf T}+4(Tc)b^{\mathsf T}.
 \label{eq:algebraic-contract-3}
\end{align}
For \eqref{eq:algebraic-contract-1}, the $(i,(k,j))$ entry of the left side is
$\sum_{k'}c_{k'}\Xi_{(k',i),(k,j)}
=\sum_{k'}c_{k'}(8T_{k'i}T_{kj}+4T_{k'k}T_{ij})
=8(Tc)_iT_{kj}+4(Tc)_kT_{ij}$, which is the corresponding entry of
$\Delta(2Tc,T)$. Multiplying this entry by $c_k\delta_{ji'}$ and summing over
$(k,j)$ gives $8(Tc)_i(Tc)_{i'}+4(c^{\mathsf T}Tc)T_{ii'}$, which is
\eqref{eq:algebraic-contract-2}. Finally $(\Delta(b,T)A_c)_{ii'}
=\sum_k(2b_kT_{ii'}+4b_iT_{ki'})c_k=2(b^{\mathsf T}c)T_{ii'}+4b_i(Tc)_{i'}$;
adding the transpose gives \eqref{eq:algebraic-contract-3}.

Let $h=g(\cdot+c)$. Then $h(x)=h_0+b'^{\mathsf T}x+x^{\mathsf T}Tx$ with
$b'=b+2Tc$ and $h_0=g(c)=g_0+b^{\mathsf T}c+c^{\mathsf T}Tc$. Write
$\Gamma(h)=\left(\begin{smallmatrix}C'&\Delta'\\\Delta'^{\mathsf T}&\Xi
\end{smallmatrix}\right)$ with $C'=2b'b'^{\mathsf T}+4h_0T$,
$\Delta'=\Delta(b',T)$, and $\Xi=\Xi(T)$. Block multiplication gives
\[
 \Psi_c^{\mathsf T}\Gamma(h)\Psi_c=
 \begin{pmatrix}
 C'-\Delta'A_c-A_c^{\mathsf T}\Delta'^{\mathsf T}+A_c^{\mathsf T}\Xi A_c
   &\Delta'-A_c^{\mathsf T}\Xi\\
 (\Delta'-A_c^{\mathsf T}\Xi)^{\mathsf T}&\Xi
 \end{pmatrix}.
\]
By \eqref{eq:algebraic-contract-1} and linearity of $\Delta$ in its first
argument, $\Delta'-A_c^{\mathsf T}\Xi=\Delta(b+2Tc,T)-\Delta(2Tc,T)=\Delta(b,T)$.
By \eqref{eq:algebraic-contract-2} and \eqref{eq:algebraic-contract-3} with
$b'$ in place of $b$, the upper left block equals
\[
 2b'b'^{\mathsf T}-4b'(Tc)^{\mathsf T}-4(Tc)b'^{\mathsf T}+8(Tc)(Tc)^{\mathsf T}
 +4\bigl(h_0-b'^{\mathsf T}c+c^{\mathsf T}Tc\bigr)T .
\]
Substituting $b'=b+2Tc$, the first four terms equal $2(b'-2Tc)(b'-2Tc)^{\mathsf
T}=2bb^{\mathsf T}$, and
$h_0-b'^{\mathsf T}c+c^{\mathsf T}Tc
=g_0+b^{\mathsf T}c+c^{\mathsf T}Tc-b^{\mathsf T}c-2c^{\mathsf T}Tc
 +c^{\mathsf T}Tc=g_0$. Hence $\Psi_c^{\mathsf T}\Gamma(h)\Psi_c=\Gamma(g)$.

\emph{(c)} The inverse of $\Psi_c$ is $\Psi_{-c}$, and
$\norm{\Psi_{-c}}\le1+\norm{c\otimes I_n}=1+\norm c$, since
$\norm{c\otimes v}=\norm c\norm v$. If $\Gamma_S(g(\cdot+c))\succeq\lambda I$
then, by (b), $y^{\mathsf T}\Gamma_S(g)y\ge\lambda\norm{\Psi_cy}^2
\ge\lambda\norm y^2/\norm{\Psi_c^{-1}}^2$ for every vector $y$.
\end{proof}

The following elementary bounds are used repeatedly.

\begin{lemma}[Block bounds]
\label{lem:algebraic-block-bounds}
For a vector $b\in\R^n$ and symmetric $T\in\R^{n\times n}$,
$\norm{\Delta(b,T)}\le\norm{\Delta(b,T)}_F\le6\norm b\norm T_F
\le6\sqrt n\norm b\norm T$ and $\Xi(T)\succeq-4\norm T^2I$. If
$T\succeq\mu I$ with $\mu\ge0$, then $\Xi(T)\succeq4\mu^2I$; if $T\succeq0$, then
$\Xi(T)\succeq0$.
\end{lemma}

\begin{proof}
The two summands of $\Delta(b,T)$ have squared Frobenius norms
$\sum_{i,k,j}4b_k^2T_{ij}^2=4\norm b^2\norm T_F^2$ and
$16\norm b^2\norm T_F^2$, and $\norm T_F\le\sqrt n\norm T$. The term
$8\operatorname{vec}(T)\operatorname{vec}(T)^{\mathsf T}$ is positive
semidefinite, and the eigenvalues of $T\otimes T$ are the products
$\lambda_i\lambda_j$ of eigenvalues of $T$.
\end{proof}

\subsection{The bivariate quartic}
\label{app:algebraic-bivariate}

\begin{proof}[Proof of Proposition~\ref{prop:algebraic-bivariate}]
Put $r=2^{1/3}$, $p=(r,r^2)$, and $q_3=(x-y)^2-2x-y+4$. Using $r^3=2$ and
$r^4=2r$ one checks $q_1(p)=q_2(p)=q_3(p)=0$. Expanding the right side
coefficient by coefficient proves the identities
\begin{equation}
\begin{gathered}
 A=7599q_1+2937q_2+5000q_3,\\
 (2r-1)q_1+(r^2-1)q_2+q_3
  =(x-r,\,y-r^2)\begin{pmatrix}2r&-1\\-1&r^2\end{pmatrix}
   (x-r,\,y-r^2)^{\mathsf T};
\end{gathered}
 \label{eq:algebraic-bivariate-expose}
\end{equation}
in the second, the constant terms agree because $r^6=4$, and the linear terms
because $r^4=2r$. In particular $A(p)=0$, so $F(p)=0$.

\emph{Zero set (b).} Since $F$ is a sum of squares, $F\ge0$ and
$F^{-1}(0)=\{F\le0\}$ is the common zero set of $A,q_1,q_2$. Now $q_1=0$
gives $y=x^2$, and then $q_2=x(x^3-2)$; the choice $x=0$ forces $y=0$ and
$A(0,0)=20000\ne0$. Hence the only common zero is $p$. The polynomial $T^3-2$
is irreducible by Eisenstein's criterion at $2$, so $r$ is irrational and
$\{F\le0\}$ contains no rational point. All coefficients of $F$ are integers;
the largest in absolute value is the coefficient $755983876$ of $x^2$, which
is below $2^{30}$.

\emph{Local data.} Put $G=A/5000$, $\varepsilon=1/2500$, and
$f=F/5000^2=G^2+\varepsilon(q_1^2+q_2^2)$. For $u\in\R^2$,
\[
 q_1(p+u)=j_1^{\mathsf T}u+u_1^2,\quad
 q_2(p+u)=j_2^{\mathsf T}u+u_2^2,\quad
 G(p+u)=\ell^{\mathsf T}u+u^{\mathsf T}Hu,
\]
where $j_1=(2r,-1)$, $j_2=(-2,2r^2)$,
$H=\left(\begin{smallmatrix}12599/5000&-1\\-1&7937/5000\end{smallmatrix}\right)$,
and, by \eqref{eq:algebraic-bivariate-expose} (whose right side has zero
gradient at $p$), $\ell=-d_1j_1-d_2j_2$ with $d_1=2r-1-\tfrac{7599}{5000}$ and
$d_2=r^2-1-\tfrac{2937}{5000}$. We claim
\begin{equation}
 H\succeq\tfrac{2937}{5000}I\succ\tfrac12I,\quad \norm H<5,\quad
 \norm{j_1},\norm{j_2}<4,\quad J^{\mathsf T}J\succeq\tfrac{36}{23}I\succ I,
 \quad \norm\ell<\tfrac1{5000},
 \label{eq:algebraic-bivariate-bounds}
\end{equation}
where $J$ has rows $j_1^{\mathsf T},j_2^{\mathsf T}$. Gershgorin's theorem
gives the first bound, and $\tr H<5$ the second. Cubing shows
$1.259921<r<1.259922$; hence $r<13/10$, $\norm{j_1}^2=4r^2+1<8$, and
$\norm{j_2}^2=4+8r<15$. Also $\det J=4r^3-2=6$ and
$\tr(J^{\mathsf T}J)=5+4r^2+8r<23$, so the smaller eigenvalue of
$J^{\mathsf T}J$ exceeds $36/23$. Finally the interval for $r$ gives
$0<d_1<4.4\cdot10^{-5}$ and $0<d_2<3.5\cdot10^{-6}$, so
$\norm\ell\le4(d_1+d_2)<4\cdot4.75\cdot10^{-5}<1/5000$.

\emph{Global Hessian bound (a).} Write $f(p+u)=f_2+f_3+f_4$ with homogeneous
parts
\[
 f_2=(\ell^{\mathsf T}u)^2+\varepsilon\norm{Ju}^2,\quad
 f_3=2(\ell^{\mathsf T}u)h(u)+2\varepsilon\bigl((j_1^{\mathsf T}u)u_1^2
     +(j_2^{\mathsf T}u)u_2^2\bigr),\quad
 f_4=h(u)^2+\varepsilon(u_1^4+u_2^4),
\]
where $h(u)=u^{\mathsf T}Hu$. Then $\nabla^2f_2\succeq2\varepsilon J^{\mathsf
T}J\succeq2\varepsilon I$. Next,
$\nabla^2(h^2)=8(Hu)(Hu)^{\mathsf T}+4h(u)H\succeq4\cdot\tfrac12\norm u^2\cdot
\tfrac12I=\norm u^2I$, and $u_1^4,u_2^4$ are convex, so
$\nabla^2f_4\succeq\norm u^2I$. For a vector $c$ and a symmetric $S$,
\[
 \nabla^2\bigl[2(c^{\mathsf T}u)(u^{\mathsf T}Su)\bigr]
  =4\bigl[(c^{\mathsf T}u)S+c(Su)^{\mathsf T}+(Su)c^{\mathsf T}\bigr],
\]
whose norm is at most $12\norm c\norm S\norm u$. Applying this to the three
terms of $f_3$, with $\norm{e_ie_i^{\mathsf T}}=1$ and
\eqref{eq:algebraic-bivariate-bounds}, gives
$\norm{\nabla^2f_3}\le(60/5000+96\varepsilon)\norm u=(63/1250)\norm u$. With
$s=\norm u$,
\[
 \nabla^2f(p+u)\succeq\Bigl(s^2-\tfrac{63}{1250}s+\tfrac1{1250}\Bigr)I
 =\Bigl(\bigl(s-\tfrac{63}{2500}\bigr)^2+\tfrac{1031}{6250000}\Bigr)I
 \succeq\tfrac{1031}{6250000}I .
\]
Multiplying by $5000^2$ gives $\nabla^2F\succeq4124I$ on $\R^2$.

\emph{The canonical Gram (c).} Order the coordinates of $w(x,v)$ as
$(v_1,v_2,x_1v_1,x_1v_2,x_2v_1,x_2v_2)$. Applying \eqref{eq:algebraic-canonical}
to the three factors $A$, $100q_1=100x^2-100y$, and $100q_2=100y^2-200x$, and
adding, gives the integer matrix
\begin{equation}
 \Gamma_F=\begin{pmatrix}
 1511967752&-6948&-1199979156&476220000&-9602&-378057752\\
 -6948&952449602&-476239204&-43876&377970000&-599989578\\
 -1199979156&-476239204&1904937612&-755940000&-755940000&899986104\\
 476220000&-43876&-755940000&599993052&300000000&-476220000\\
 -9602&377970000&-755940000&300000000&599993052&-476220000\\
 -378057752&-599989578&899986104&-476220000&-476220000&756071628
 \end{pmatrix}.
 \label{eq:algebraic-bivariate-gram}
\end{equation}
For instance, its first diagonal entry is $2\cdot15874^2+4\cdot20000\cdot12599
+2\cdot200^2$. By Lemma~\ref{lem:algebraic-gram}(a),
$v^{\mathsf T}\nabla^2F(x)v=w(x,v)^{\mathsf T}\Gamma_Fw(x,v)$.

To prove $\Gamma_F\succ0$, recentre the three factors at $p$. Their recentred
canonical Gram is $5000^2M$ with
$M=\left(\begin{smallmatrix}C&\Delta\\\Delta^{\mathsf T}&\Xi\end{smallmatrix}
\right)$, where, because the recentred factors vanish at the origin,
\[
 C=2\ell\ell^{\mathsf T}+2\varepsilon J^{\mathsf T}J,\quad
 \Delta=\Delta(\ell,H)+\varepsilon\bigl(\Delta(j_1,e_1e_1^{\mathsf T})
          +\Delta(j_2,e_2e_2^{\mathsf T})\bigr),\quad
 \Xi=\Xi(H)+\varepsilon\bigl(\Xi(e_1e_1^{\mathsf T})+\Xi(e_2e_2^{\mathsf T})\bigr).
\]
By \eqref{eq:algebraic-bivariate-bounds} and
Lemma~\ref{lem:algebraic-block-bounds}, $C\succeq\tfrac{72}{23}\varepsilon I
\succeq0.00125\,I$ and $\Xi\succeq4(2937/5000)^2I\succeq1.38\,I$. Since
$\norm H_F^2<11$,
$\norm\Delta\le6\bigl(\norm\ell\norm H_F+\varepsilon(\norm{j_1}+\norm{j_2})\bigr)
<6(3.32/5000+8/2500)<0.0232$. The Schur complement therefore satisfies
$C-\Delta\Xi^{-1}\Delta^{\mathsf T}\succeq(0.00125-0.0232^2/1.38)I\succ0$, so
$M\succ0$. Lemma~\ref{lem:algebraic-gram}(b) gives
$\Gamma_F=5000^2\Psi_p^{\mathsf T}M\Psi_p\succ0$.

\emph{A direct rational certificate (d).} For $x=(x_1,x_2)$ and
$v=(a,b)$ put $X=100x_1-126$, $Y=100x_2-\tfrac{1587}{10}$, and
\[
 \omega=\bigl(2a-2b,\ 2b,\ 2Xa-Xb-Ya+Yb,\ 2Xb+Ya-Yb,\ 2Ya-Yb,\ 2Yb\bigr)^{\mathsf T}.
\]
Let $N_F$ be the symmetric integer matrix
\begin{center}
\footnotesize
\setlength{\arraycolsep}{2.5pt}
$N_F=\begin{pmatrix}
 812297303152&102427849072&22618204800&-13093175680&5654551200&-130053040\\
 102427849072&262472489120&22618204800&-5643994880&1929960800&-3911932400\\
 22618204800&22618204800&761975044800&78611522400&39305761200&-2034439000\\
 -13093175680&-5643994880&78611522400&128114982000&-15939729800&-13708914300\\
 5654551200&1929960800&39305761200&-15939729800&112025966300&13144848050\\
 -130053040&-3911932400&-2034439000&-13708914300&13144848050&120298532375
 \end{pmatrix}$.
\end{center}
Expanding both sides, using $v^{\mathsf T}\nabla^2F(x)v=w^{\mathsf
T}\Gamma_Fw$ and \eqref{eq:algebraic-bivariate-gram}, proves the polynomial
identity
\begin{equation}
 16\cdot10^6\bigl(v^{\mathsf T}\nabla^2F(x)v-4096\norm v^2\bigr)
 =\omega^{\mathsf T}N_F\,\omega .
 \label{eq:algebraic-bivariate-squares}
\end{equation}
Every row of $N_F$ is strictly diagonally dominant; the margins
$\delta_i=(N_F)_{ii}-\sum_{j\ne i}\abs{(N_F)_{ij}}$ are
\[
\begin{aligned}
 (\delta_0,\delta_1,\delta_2)&=(668373469360,\ 125940547168,\ 596786912600),\\
 (\delta_3,\delta_4,\delta_5)&=(1117644940,\ 36051115250,\ 87368345585).
\end{aligned}
\]
Hence
$\omega^{\mathsf T}N_F\omega=\sum_i\delta_i\omega_i^2+\sum_{i<j}\abs{(N_F)_{ij}}
\bigl(\omega_i+\operatorname{sign}((N_F)_{ij})\,\omega_j\bigr)^2$, a positive
integer combination of $21$ squares of polynomials with rational coefficients.
This proves $\nabla^2F\succeq4096I$ by a rational certificate; equivalently,
$F-2048\norm x^2$ is SOS-convex with an explicit rational Hessian sum of
squares.
\end{proof}

\subsection{One real conjugate}
\label{app:algebraic-conjugate}

We use two standard facts. First, the field $\R_{\mathrm{alg}}$ of real
algebraic numbers is real closed, and by the Tarski--Seidenberg transfer
principle a system of polynomial equations and inequalities with rational
coefficients that has a solution in $\R^n$ has a solution in
$\R_{\mathrm{alg}}^n$ \cite[Theorem~2.80 and Corollary~2.83]{BasuPollackRoy2006}.
Second, a number field $K$ of degree $[K:\Q]$ has $r_1$ real embeddings and
$2r_2$ nonreal embeddings occurring in conjugate pairs, with
$[K:\Q]=r_1+2r_2$.

\begin{proof}[Proof of Lemma~\ref{lem:one-real-conjugate}]
\emph{Algebraicity.} In both cases the singleton is the solution set of a
finite system of rational polynomial equations or weak inequalities. The
system has a real solution, hence a solution in $\R_{\mathrm{alg}}^n$, which
is also a real solution and therefore equals $p$.

\emph{Reduction to equalities.} In the zero-set case let $E$ be the given
rational polynomials. In the convex case let
$\mathcal A=\{i:f_i(p)=0\}$ and $E=\{f_i:i\in\mathcal A\}$. We claim that
$\{x:f_i(x)\le0,\ i\in\mathcal A\}=\{p\}$. Suppose $z\ne p$ satisfies these
inequalities. For $0\le s\le1$ and $i\in\mathcal A$, convexity gives
$f_i(p+s(z-p))\le(1-s)f_i(p)+sf_i(z)\le0$. For $i\notin\mathcal A$ we have
$f_i(p)<0$, so by continuity and finiteness there is $s_0>0$ with
$f_i(p+s(z-p))<0$ for all $i\notin\mathcal A$ and $0\le s\le s_0$. Then
$p+s_0(z-p)\ne p$ satisfies all original inequalities, a contradiction.

\emph{One real embedding.} Put $K=\Q(p)$ and let $\sigma:K\to\R$ be a field
embedding. Apply $\sigma$ coordinatewise to $p$. For every $e\in E$, rational
coefficients give $e(\sigma(p))=\sigma(e(p))=0$. In the zero-set case
$\sigma(p)$ is therefore a real common zero, so $\sigma(p)=p$. In the convex
case $\sigma(p)$ satisfies $f_i(\sigma(p))=0\le0$ for $i\in\mathcal A$, so
$\sigma(p)=p$ by the claim. As $K$ is generated by the coordinates of $p$,
$\sigma$ is the inclusion. Hence $r_1=1$, and $[K:\Q]=1+2r_2$ is odd.

\emph{Coordinates.} Let $K_j=\Q(p_j)\subseteq K$ and let $\tau:K_j\to\R$ be
a real embedding. Since $[K:\Q]$ is odd, so is $[K:K_j]$. Choose a primitive
element $\vartheta$ of $K/K_j$ with minimal polynomial $m_\vartheta\in K_j[T]$
of odd degree $[K:K_j]$. The real polynomial $\tau(m_\vartheta)$ has odd
degree, hence a real root $\vartheta'$, and $\vartheta\mapsto\vartheta'$
extends $\tau$ to a real embedding of $K=K_j(\vartheta)$. Distinct real
embeddings of $K_j$ would give distinct real embeddings of $K$. So $K_j$ has
exactly one real embedding, which means that the minimal polynomial of $p_j$
has exactly one real root.
\end{proof}

\begin{proof}[Proof of Theorem~\ref{thm:singleton-field}]
(v)$\Rightarrow$(ii): the real zero set of $f$ is the given singleton.
(v)$\Rightarrow$(iii): $f$ is a sum of squares, so $f\ge0$ and
$\{f\le0\}=f^{-1}(0)$; take the single convex inequality $f\le0$. (v)$\Rightarrow$(iv): the unique zero of
the nonnegative strongly convex $f$ is its unique minimizer.

(ii)$\Rightarrow$(i) and (iii)$\Rightarrow$(i) are
Lemma~\ref{lem:one-real-conjugate}. For (iv)$\Rightarrow$(i), let $p$ be the
unique minimizer of a rational convex polynomial $h$. For a convex
differentiable function the critical points are exactly the global
minimizers, so the real zero set of the rational polynomial
$\sum_i(\partial_ih)^2$ is $\{p\}$, and Lemma~\ref{lem:one-real-conjugate}
applies.

(i)$\Rightarrow$(v): if $\alpha$ is rational, $f=(x-\alpha)^2+(x-\alpha)^4$
is a sum of two rational squares with $f''=2+12(x-\alpha)^2\ge2$. Writing
$f$ as $(x-\alpha)^2+(x^2-2\alpha x+\alpha^2)^2$, its canonical Hessian Gram
on $w=(v,xv)$ is $\left(\begin{smallmatrix}2+12\alpha^2&-12\alpha\\-12\alpha&12
\end{smallmatrix}\right)$, which has determinant $24$ and is positive
definite. If $\alpha$ is irrational, its minimal polynomial has degree $d>1$
and exactly one real root, so the nonreal roots pair up and $d\ge3$ is odd.
Theorem~\ref{thm:algebraic-dimension} with $n=(d+1)/2$ gives (v), with
$\nabla^2f\succeq\tfrac32I\succeq I$ and $n+1$ rational squares. Its
algorithm accepts any rational multiple of the minimal polynomial, because the
first step normalizes the input to a primitive integer polynomial.
\end{proof}


\subsection{The generic realization lemma}
\label{app:algebraic-realization}

\begin{proof}[Proof of Lemma~\ref{lem:quartic-realization}]
Put $\Phi=g^2+\varepsilon\sum_jr_j^2=\varepsilon\nu^2f$. The representation
of $f$ uses the $m+1$ rational squares $g/(t\nu)$ and $r_j/\nu$, and every
factor vanishes at $p$.

\emph{Centred Gram.} Let $M$ be the canonical Hessian Gram of the recentred
representation $\Phi(p+\cdot)=g(p+\cdot)^2+\varepsilon\sum_jr_j(p+\cdot)^2$.
The recentred factors have no constant terms, so by
\eqref{eq:algebraic-canonical}
\[
\begin{gathered}
 M=\begin{pmatrix}C&\Delta\\\Delta^{\mathsf T}&\Xi\end{pmatrix},\qquad
 C=2\ell\ell^{\mathsf T}+2\varepsilon\sum_jb_jb_j^{\mathsf T},\\
 \Delta=\Delta(\ell,H)+\varepsilon\sum_j\Delta(b_j,T_j),\qquad
 \Xi=\Xi(H)+\varepsilon\sum_j\Xi(T_j).
\end{gathered}
\]
By \eqref{eq:algebraic-realization-data} and
Lemma~\ref{lem:algebraic-block-bounds}, $C\succeq2\varepsilon\nu^2I$,
$\Xi\succeq(4\mu^2-4\varepsilon m)I\succeq2\mu^2I$ (using
$\varepsilon\le\mu^2/(2m)$), and
\[
 \norm\Delta\le6\sqrt n\bigl(\norm\ell\Lambda+\varepsilon m\beta\bigr)
 \le6\sqrt n\,\varepsilon(\Lambda+m\beta)=:B_\Delta .
\]
The individual matrices $\Xi(T_j)$ may be indefinite; the bound
$\Xi(T_j)\succeq-4I$ is what the estimate uses. The Schur complement satisfies
\[
 S:=C-\Delta\Xi^{-1}\Delta^{\mathsf T}
 \succeq\Bigl(2\varepsilon\nu^2-\frac{B_\Delta^2}{2\mu^2}\Bigr)I
 =\Bigl(2\varepsilon\nu^2-\frac{18n\varepsilon^2(\Lambda+m\beta)^2}{\mu^2}\Bigr)I
 \succeq\tfrac32\varepsilon\nu^2I,
\]
where the last step uses $\varepsilon\le\nu^2\mu^2/(36n(\Lambda+m\beta)^2)$.
For $y=(y_1,y_2)$ conformal with the blocks, put
$y_2'=y_2+\Xi^{-1}\Delta^{\mathsf T}y_1$. Completing the square,
\begin{equation}
 y^{\mathsf T}My=y_1^{\mathsf T}Sy_1+y_2'^{\mathsf T}\Xi y_2'.
 \label{eq:algebraic-schur}
\end{equation}

\emph{(b).} Lemma~\ref{lem:algebraic-gram}(a), applied to the recentred
representation, gives
\[
 v^{\mathsf T}\nabla^2\Phi(p+u)\,v=w(u,v)^{\mathsf T}Mw(u,v)
 \qquad(u,v\in\R^n).
\]
The first block of $w(u,v)$ is $v$, so \eqref{eq:algebraic-schur} gives
$v^{\mathsf T}\nabla^2\Phi(p+u)v\ge\tfrac32\varepsilon\nu^2\norm v^2$ for all
$u,v$. Dividing by $\varepsilon\nu^2$ gives $\nabla^2f\succeq\frac32I$. A
nonnegative strongly convex function vanishing at $p$ has $p$ as its unique
zero and unique minimizer.

\emph{(a).} The quartic part of $\Phi$ is
$(x^{\mathsf T}Hx)^2+\varepsilon\sum_j(x^{\mathsf T}T_jx)^2\ge\mu^2\norm x^4$,
so $f$ has degree exactly four and positive definite quartic form.

\emph{(c).} With $\theta'=\norm{\Xi^{-1}\Delta^{\mathsf T}}\le
B_\Delta/(2\mu^2)=\theta$, we have
$\norm y^2\le\bigl((1+\theta)\norm{y_1}+\norm{y_2'}\bigr)^2
\le((1+\theta)^2+1)(\norm{y_1}^2+\norm{y_2'}^2)$. Hence
\eqref{eq:algebraic-schur} yields
$M\succeq\min\{\tfrac32\varepsilon\nu^2,2\mu^2\}((1+\theta)^2+1)^{-1}I$. The
canonical Gram of the representation of $\Phi$ is $\varepsilon\nu^2\Gamma_f$;
by Lemma~\ref{lem:algebraic-gram}(b) and (c) with $c=p$,
$\varepsilon\nu^2\Gamma_f=\Psi_p^{\mathsf T}M\Psi_p\succeq
\lambda_{\min}(M)(1+\norm p)^{-2}I$. Dividing by $\varepsilon\nu^2$ gives
$\Gamma_f\succeq\gamma I$. By Lemma~\ref{lem:algebraic-gram}(a), $\Gamma_f$
is a full Hessian Gram of $f$, and it is rational because it is computed by
\eqref{eq:algebraic-canonical} from the rational factors. Each of the $m+1$
matrices $\Gamma(\cdot)$ has $(n+n^2)^2$ entries, each a fixed polynomial of
degree two in the coefficients.
\end{proof}

\begin{remark}[Curvature without the factor $n$]
\label{rem:algebraic-curvature}
Part (b) of Lemma~\ref{lem:quartic-realization} remains true when the last
bound in \eqref{eq:algebraic-realization-eps} is replaced by
$\varepsilon\le\nu^2\mu^2/(36(\Lambda+m\beta)^2)$, the other conditions being
kept. In the proof above the factor $n$ enters only through the operator-norm
bound $\norm{\Delta(b,T)}\le6\sqrt n\norm b\norm T$, while part (b) needs the
Gram only at vectors of the form $w(u,v)$. A direct estimate at these vectors
avoids it. For a quadratic $q(p+u)=c^{\mathsf T}u+u^{\mathsf T}Su$,
\[
 v^{\mathsf T}\nabla^2(q^2)(p+u)\,v
 =2\bigl(c^{\mathsf T}v+2u^{\mathsf T}Sv\bigr)^2
 +4\bigl(c^{\mathsf T}u+u^{\mathsf T}Su\bigr)\,v^{\mathsf T}Sv .
\]
Fix $u,v$ and put $s=\norm u$. The terms of this expression of degree zero
in $u$ are $2(c^{\mathsf T}v)^2$; those of degree one,
$8(c^{\mathsf T}v)(u^{\mathsf T}Sv)+4(c^{\mathsf T}u)(v^{\mathsf T}Sv)$, have
absolute value at most $12\norm c\norm S\,s\norm v^2$; and those of degree
two are $8(u^{\mathsf T}Sv)^2+4(u^{\mathsf T}Su)(v^{\mathsf T}Sv)$. Apply this
to $g$ ($c=\ell$, $S=H$) and to each $r_j$ ($c=b_j$, $S=T_j$), and add with the
weights of $\Phi=g^2+\varepsilon\sum_jr_j^2$. By
\eqref{eq:algebraic-realization-data}, the terms of degree zero sum to at
least $2\varepsilon\nu^2\norm v^2$. The terms of degree two are at least
$4\mu^2s^2\norm v^2$ for $g$ and at least $-4s^2\norm v^2$ for each $r_j$, so
their weighted sum is at least $(4\mu^2-4\varepsilon m)s^2\norm v^2\ge
2\mu^2s^2\norm v^2$. The terms of degree one have weighted sum of absolute
value at most $12(\norm\ell\Lambda+\varepsilon m\beta)s\norm v^2\le
12\varepsilon(\Lambda+m\beta)s\norm v^2$. Hence
\[
 v^{\mathsf T}\nabla^2\Phi(p+u)\,v\ge
 \bigl(2\varepsilon\nu^2-12\varepsilon(\Lambda+m\beta)s+2\mu^2s^2\bigr)
 \norm v^2 .
\]
The minimum of the bracket over $s\ge0$ is
$2\varepsilon\nu^2-18\varepsilon^2(\Lambda+m\beta)^2/\mu^2$, which is at least
$\frac32\varepsilon\nu^2$ under the stated bound on $\varepsilon$. Dividing by
$\varepsilon\nu^2$ gives $\nabla^2f\succeq\frac32I$. The proof of part (c)
keeps the factor $n$, because its matrix inequality concerns all vectors of
$\R^{n+n^2}$, not only those of the form $w(u,v)$.
\end{remark}

A rational positive semidefinite Gram matrix also yields an explicit rational
sum of squares of the Hessian biform. Rational symmetric elimination writes
$\Gamma_f=\sum_kc_k\lambda_k\lambda_k^{\mathsf T}$ with rational vectors
$\lambda_k$ and rational $c_k>0$. A positive rational $u/v$, with $u,v$
positive integers, equals $uv/v^2$; writing $uv=\sum_h\epsilon_h2^h$ in binary
and using $2^{2h}/v^2=(2^h/v)^2$ and $2^{2h+1}/v^2=2(2^h/v)^2$ expresses it as
a sum of at most $2\lceil\log_2(uv+1)\rceil$ rational squares, without
factoring integers. This conversion is not needed for any result of
Section~\ref{sec:algebraic}.

\subsection{Power coordinates}
\label{app:algebraic-power}

The algorithm uses one imported numerical subroutine.

\begin{lemma}[Certified root approximation; imported]
\label{lem:algebraic-roots}
There is a deterministic algorithm that, given a square-free $P\in\Z[T]$ of
degree $d$ with coefficients of absolute value at most $2^\tau$ and an integer
$N_0\ge1$, outputs Gaussian dyadic numbers $z_1,\ldots,z_d$ and the guarantee
that the complex roots of $P$ can be enumerated as $\beta_1,\ldots,\beta_d$
with $\abs{z_i-\beta_i}\le2^{-N_0}$ for every $i$. Its bit complexity is
polynomial in $d$, $\tau$, and $N_0$.
\end{lemma}

This is a classical result; one primary source is
\cite[Theorem~5]{MehlhornSagraloffWang2015}, which isolates all complex roots
and refines the isolating disks to any requested radius in
$\widetilde O(d^3+d^2\tau+dN_0)$ bit operations; see also
\cite{Schonhage1982,Pan2002}. Only polynomiality is used. An uncertified
numerical root finder cannot replace this lemma.

\begin{proof}[Proof of Theorem~\ref{thm:algebraic-dimension}]
\emph{Normalization and root bounds.} Clearing denominators and dividing by
the content and the sign of the leading coefficient gives a primitive
$P_{\Z}=\sum_{i=0}^da_iT^i\in\Z[T]$ with $a_d>0$ and $\abs{a_i}\le2^\tau$,
$\tau\ge1$, in polynomial time; $\tau$ is polynomial in $L$. Let
$P_1=P_{\Z}/a_d=\sum_ic_iT^i$ be the monic minimal polynomial of $\alpha$, so
$\abs{c_i}\le2^\tau$. Put
\[
 R=2^{\tau+1},\qquad \sigma=2^{-\tau(d-1)}(2R)^{-d(d-1)/2}.
\]
Every complex root $\beta$ satisfies $\abs\beta\le1+\max_i\abs{c_i}<R$. The
discriminant of $P_{\Z}$ is a nonzero integer, since an irreducible polynomial
over $\Q$ is separable, and equals $a_d^{2d-2}\prod_{i<j}(\beta_i-\beta_j)^2$.
Bounding all but one difference by $2R$ shows that distinct roots satisfy
$\abs{\beta-\beta'}\ge a_d^{-(d-1)}(2R)^{-(d(d-1)/2-1)}>\sigma$. Hence every
nonreal root has $\abs{\operatorname{Im}\beta}>\sigma/2$, and
$\abs{P_1'(\alpha)}=\prod_{\beta\ne\alpha}\abs{\alpha-\beta}>\sigma^{d-1}$.
The nonreal roots form $s=(d-1)/2$ conjugate pairs; let
$\beta_j=u_j+\mathrm iv_j$, $v_j>\sigma/2$, $1\le j\le s$, be those in the
upper half-plane. Write $v_k(T)=(1,T,\ldots,T^k)^{\mathsf T}$.

\emph{A positive definite Gram of $P_1(T)(1+T^2)^e/(T-\alpha)$.} Put
$e=n-s-1\ge0$. For $1\le j\le s$ let
$B_j=\left(\begin{smallmatrix}u_j^2+v_j^2&-u_j\\-u_j&1\end{smallmatrix}\right)$,
so that $v_1^{\mathsf T}B_jv_1=(T-\beta_j)(T-\overline{\beta_j})$, and for
$s<j\le n-1$ let $B_j=I_2$, so that $v_1^{\mathsf T}B_jv_1=1+T^2$. Since
$\det B_j=v_j^2>\sigma^2/4$ and $\tr B_j<1+R^2$,
\[
 b_0I\preceq B_j,\qquad \norm{B_j}\le1+R^2,\qquad
 b_0=\frac{\sigma^2}{4(1+R^2)} .
\]
Let $S_j$ be the $0/1$ matrix with
$S_jv_j(T)=(v_{j-1}(T)^{\mathsf T},Tv_{j-1}(T)^{\mathsf T})^{\mathsf T}$; then
$S_j^{\mathsf T}S_j$ is diagonal with entries $1$ and $2$, so
$I\preceq S_j^{\mathsf T}S_j\preceq2I$. Define $G_0=[1]$ and
$G_j=S_j^{\mathsf T}(B_j\otimes G_{j-1})S_j$. Because
$S_jv_j=v_1\otimes v_{j-1}$, the mixed-product rule gives
$v_j^{\mathsf T}G_jv_j=(v_1^{\mathsf T}B_jv_1)(v_{j-1}^{\mathsf
T}G_{j-1}v_{j-1})$. Hence $G'=G_{n-1}$ satisfies
\[
 v_{n-1}^{\mathsf T}G'v_{n-1}=\frac{P_1(T)}{T-\alpha}(1+T^2)^e,\qquad
 G'\succeq b_0^{n-1}I,\qquad \norm{G'}\le\bigl(2(1+R^2)\bigr)^{n-1}.
\]

\emph{The exposing matrix.} Let $D_c$ be the $n\times(n+1)$ matrix with
$(D_cz)_k=z_{k+1}-cz_k$, $0\le k<n$, so that $D_cv_n(T)=(T-c)v_{n-1}(T)$, and
put $Q_*=D_\alpha^{\mathsf T}G'D_\alpha$. Then
\begin{equation}
 v_n(T)^{\mathsf T}Q_*v_n(T)=(T-\alpha)P_1(T)(1+T^2)^e,\qquad Q_*\succeq0,
 \qquad \ker Q_*=\ker D_\alpha=\R v_n(\alpha).
 \label{eq:algebraic-qstar}
\end{equation}
For $y\in\R^n$, the equation $D_\alpha(0,y)=\omega$ is solved by
$y_k=\sum_{i<k}\alpha^{k-1-i}\omega_i$, so $\norm y\le nR^{n-1}\norm\omega$,
and therefore
\begin{equation}
 (0,y)^{\mathsf T}Q_*(0,y)\ge\gamma_0\norm y^2,\quad
 \gamma_0=\frac{b_0^{n-1}}{n^2R^{2n-2}},\qquad
 \norm{Q_*}\le\Lambda_*=(1+R)^2\bigl(2(1+R^2)\bigr)^{n-1}.
 \label{eq:algebraic-qstar-bounds}
\end{equation}

\emph{The rational vanishing space.} Let $\mathcal L_n$ be the set of
symmetric $(n+1)\times(n+1)$ matrices $Q$ with
$v_n(T)^{\mathsf T}Qv_n(T)\equiv0\pmod{P_1}$. For $0\le k<d$ let $E_k$ be the
rational symmetric matrix whose $(i,j)$ entry is the coefficient of $T^k$ in
the remainder of $T^{i+j}$ modulo $P_1$. With $\ip XY=\tr(XY)$,
$v_n^{\mathsf T}Qv_n\equiv\sum_k\ip{E_k}QT^k$, so $\mathcal L_n$ is the common
kernel of the forms $\ip{E_k}{\cdot}$, and $Q_*\in\mathcal L_n$. The $E_k$
are linearly independent: since $d-1\le2n$, every $k<d$ equals $i+j$ for some
$0\le i,j\le n$, and at that entry $(E_l)_{ij}=\delta_{kl}$. Hence the Gram
matrix $\Gamma_E=(\ip{E_k}{E_l})$ is invertible, and
\[
 \Pi(X)=X-\sum_{k,l}(\Gamma_E^{-1})_{kl}\ip{E_l}XE_k
\]
is the orthogonal projection onto $\mathcal L_n$. It maps rational matrices
to rational matrices, and it does not increase Frobenius distances to points
of $\mathcal L_n$. For $d\le m\le2n$ write the remainder $T^m\bmod P_1$ as
$\sum_i\rho_{m,i}T^i$. Then $\rho_{m+1,i}=\rho_{m,i-1}-\rho_{m,d-1}c_i$, so
all $\abs{\rho_{m,i}}\le R^{m-d+1}$, and $a_d^{m-d+1}\rho_{m,i}$ is an
integer. All entries of $E_k$, $\Gamma_E$, and, by Cramer's rule and
Hadamard's inequality, of $\Gamma_E^{-1}$ have bit length polynomial in $d$
and $\tau$.

\emph{Approximation.} Let $0<\zeta\le\min\{1,\sigma/16\}$, and apply
Lemma~\ref{lem:algebraic-roots} with $2^{-N_0}\le\zeta$. An approximation
$z$ of the real root has $\abs{\operatorname{Im}z}\le\zeta<\sigma/4$, and an
approximation of a nonreal root has $\abs{\operatorname{Im}z}>\sigma/2-\zeta
>\sigma/4$; this identifies $\tilde\alpha=\operatorname{Re}z$ for the real
root and approximations $\tilde\beta_j$ of the upper roots, each within
$\zeta$. Form $\tilde B_j$ from $\tilde\beta_j$ as above (and $\tilde B_j=I_2$
for $j>s$), the recursion $\tilde G_j$, and the rational matrix
$\tilde X=D_{\tilde\alpha}^{\mathsf T}\tilde G_{n-1}D_{\tilde\alpha}$. With
$\hat B=1+(R+1)^2$ and $C_B=2R+3$ we have
$\norm{\tilde B_j-B_j}\le\bigl(\abs{\tilde\beta_j}+\abs{\beta_j}\bigr)\zeta
+\sqrt2\zeta\le C_B\zeta$ and $\norm{\tilde B_j},\norm{B_j}\le\hat B$. By
induction on $j$, using $\norm{S_j}^2\le2$,
\[
 \norm{\tilde G_j-G_j}\le2\bigl(C_B\zeta(2\hat B)^{j-1}
   +\hat B\norm{\tilde G_{j-1}-G_{j-1}}\bigr)\le2jC_B\zeta(2\hat B)^{j-1}.
\]
Since $\norm{D_{\tilde\alpha}-D_\alpha}\le\zeta$, $\norm{D_{\tilde\alpha}}\le
R+2$, and $\norm{D_\alpha}\le R+1$,
\[
 \norm{\tilde X-Q_*}_F\le(n+1)C_X\zeta,\qquad
 C_X=(2\hat B)^{n-1}(2R+3)\bigl((n-1)(R+2)^2+1\bigr).
\]
For a rational $0<\delta\le1$ choose $\zeta\le\delta/(2(n+1)C_X)$ as well, and
put $Q=\Pi(\tilde X)$. Then $Q\in\mathcal L_n$ is rational and
$\norm{Q-Q_*}_F\le\delta/2$.

\emph{The exposing quadratic.} Let $g(x)=(1,x)^{\mathsf T}Q(1,x)$ and
$p=(\alpha,\ldots,\alpha^n)$, so $(1,p)=v_n(\alpha)$. Since
$Q\in\mathcal L_n$, $g(p)=v_n(\alpha)^{\mathsf T}Qv_n(\alpha)=0$. For
$u\in\R^n$, $g(p+u)=\ell^{\mathsf T}u+u^{\mathsf T}Hu$ with
$\ell=2\bigl[Qv_n(\alpha)\bigr]_{1:n}=2\bigl[(Q-Q_*)v_n(\alpha)\bigr]_{1:n}$ and
$H=Q_{1:n,1:n}$. With $\kappa_v=(n+1)R^n\ge\norm{v_n(\alpha)}$,
\eqref{eq:algebraic-qstar-bounds} gives
\begin{equation}
 \norm\ell\le\kappa_v\delta,\qquad H\succeq(\gamma_0-\delta/2)I,\qquad
 \norm H\le\Lambda_*+\delta/2.
 \label{eq:algebraic-exposing-bounds}
\end{equation}

\emph{Residuals.} Put $x_0=1$; for $0\le k\le d$ let $\phi_k=x_k$ if $k\le n$
and $\phi_k=x_nx_{k-n}$ if $n<k\le d$ (then $1\le k-n\le n-1$, because
$d\le2n-1$). Let $\kappa=1+\sum_{k=n+1}^d\abs{c_k}$ and
\[
 r_j=\frac{x_1x_j-x_{j+1}}{\kappa}\ (1\le j<n),\qquad
 r_n=\frac1\kappa\sum_{k=0}^dc_k\phi_k .
\]
On the moment curve $\gamma(T)=(T,T^2,\ldots,T^n)$ we have $\phi_k(\gamma(T))
=T^k$, so $r_j(\gamma(T))=0$ for $j<n$ and $r_n(\gamma(T))=P_1(T)/\kappa$. In
particular every $r_j$ vanishes at $p=\gamma(\alpha)$. Conversely, the first
$n-1$ equations force $x_j=x_1^j$, and then $r_n=P_1(x_1)/\kappa$, so $p$ is the
only real common zero. Each $r_j$ with $j<n$ has quadratic part
$x_1x_j/\kappa$, of matrix norm at most $1$; the quadratic part of $r_n$ is a
combination of the $x_nx_{k-n}$ with coefficient sum at most $\kappa-1$.
Hence $\norm{T_j}\le1$ for all $j$. The gradients satisfy
$\norm{\nabla r_j(p)}\le3R^n$ for $j<n$ and
$\norm{\nabla r_n(p)}\le\sum_k\abs{c_k}\,2R^n\le(d+1)R^{n+1}$, so the
Jacobian $J$ at $p$ satisfies $\norm J\le\norm J_F\le\beta=(d+4)R^{n+1}$ and
every row has norm at most $\beta$.

The determinant of $J$ is $P_1'(\alpha)/\kappa^n$. Indeed, differentiating
$r(\gamma(T))=(0,\ldots,0,P_1(T)/\kappa)$ at $T=\alpha$ gives
$J\omega=(0,\ldots,0,P_1'(\alpha)/\kappa)^{\mathsf T}$ with
$\omega=\gamma'(\alpha)=(1,2\alpha,\ldots,n\alpha^{n-1})$. Replace the first
column of the identity by $\omega$; this change of basis has determinant
$\omega_1=1$. Expanding along the new first column, $\det J=(-1)^{n+1}
\kappa^{-1}P_1'(\alpha)\det J'$, where $J'$ consists of rows $1,\ldots,n-1$ and
columns $2,\ldots,n$ of $J$. Row $j$ of $J'$ has entry $-1/\kappa$ in column
$j+1$ and $\alpha/\kappa$ in column $j$ (for $j\ge2$), so $J'$ is triangular
with diagonal $-1/\kappa$ and $\det J'=(-1/\kappa)^{n-1}$. The product of the
singular values now gives
\[
 \sigma_{\min}(J)\ge\frac{\abs{\det J}}{\norm J^{n-1}}\ge
 \nu=\frac{\sigma^{d-1}}{\kappa^n\beta^{n-1}},\qquad
 \sum_jb_jb_j^{\mathsf T}=J^{\mathsf T}J\succeq\nu^2I .
\]

\emph{Parameters and output.} Put $\mu=\gamma_0/2$ and $\Lambda=\Lambda_*+1$.
Let $t=2^{-N_t}$ with $N_t$ the least nonnegative integer such that
$\varepsilon=4^{-N_t}\le\min\{1,\mu^2/(2n),\nu^2\mu^2/(36n(\Lambda+n\beta)^2)\}$;
then take $\delta=\min\{1,\gamma_0,\varepsilon/\kappa_v\}$ and the corresponding
$\zeta$. By \eqref{eq:algebraic-exposing-bounds}, $H\succeq\mu I$,
$\norm H\le\Lambda$, and $\norm\ell\le\varepsilon$. With $m=n$ residuals all
hypotheses of Lemma~\ref{lem:quartic-realization} hold, which gives every
stated property of $F_P$.

\emph{Complexity.} Every bound above ($R$, $\sigma$, $b_0$, $\gamma_0$,
$\Lambda_*$, $\kappa$, $\kappa_v$, $\beta$, $\nu$, $C_X$) is a rational number
obtained from $d$, $n$, $\tau$ and the $c_i$ by a polynomial number of
arithmetic operations, and the logarithms of these numbers and of their
reciprocals are $O(d^3\tau+d\log d)$. Thus
$N_0$, $N_t$, and $\log(1/\delta)$ are polynomial in $d$ and $\tau$. Root
approximation is polynomial by Lemma~\ref{lem:algebraic-roots}. The matrices
$\tilde B_j$, $\tilde G_j$, $D_{\tilde\alpha}$, $\tilde X$ and $Q$ have order at
most $2n$ and are obtained by $O(n^4+d^4)$ rational operations; the entries of
$\tilde G_j$ have bit length $O(j(N_0+\tau))$, and the projection uses the
polynomial-size data of $\Gamma_E^{-1}$. The output $g,r_j,t,\nu$ and the
canonical Gram therefore have polynomial bit length and are computed in
polynomial time.

\emph{Converse.} Let $f$ be as in the converse statement and
$h(T)=f(T,T^2,\ldots,T^k)\in\Q[T]$. If $h=0$, then $f$ vanishes on the whole
real moment curve, contradicting uniqueness of the zero. So $h$ is a nonzero
polynomial of degree at most $4k$, nonnegative on $\R$, with $h(\alpha)=0$.
A real zero of a nonnegative polynomial has even multiplicity, so $\alpha$ is
at least a double root. Since $P_1$ is the minimal polynomial of $\alpha$ and
is separable, $P_1\mid h$; writing $h=P_1h_1$ and using $P_1'(\alpha)\ne0$, also
$h_1(\alpha)=0$ and $P_1\mid h_1$. Hence $2d\le\deg h\le4k$, so
$k\ge d/2$, that is $k\ge(d+1)/2$.
\end{proof}

\begin{proof}[Proof of Corollary~\ref{cor:algebraic-three-quadrics}]
Use the preceding construction with $n=d-1$; write $v=v_{d-1}$ and
$p'=(\alpha,\ldots,\alpha^{d-1})$. Let $C$ be the companion matrix of $P_1$,
with $(Cz)_k=z_{k+1}$ for $k\le d-2$ and $(Cz)_{d-1}=-\sum_ic_iz_i$. Then
$Cv(\beta)=\beta v(\beta)$ for every root $\beta$. Let
$\ell_\alpha=(\ell_0,\ldots,\ell_{d-1})$ be the coefficient vector of
$P_1(T)/(T-\alpha)=\sum_i\ell_iT^i$; comparing coefficients in
$(T-\alpha)\sum_i\ell_iT^i=P_1(T)$ gives $\ell_\alpha^{\mathsf
T}C=\alpha\ell_\alpha^{\mathsf T}$. Since $\alpha$ is a simple eigenvalue,
$U:=\operatorname{image}(C-\alpha I)$ has dimension $d-1$ and equals
$\ker\ell_\alpha^{\mathsf T}$. Also $\ell_\alpha^{\mathsf
T}v(\alpha)=P_1'(\alpha)$, and $\abs{\ell_i}\le\sum_{k>i}\abs{c_k}R^{k-i-1}\le
dR^d$, so $\norm{\ell_\alpha}\le d^2R^d$.

\emph{Margin of $Q_*$ on $U$.} Let $w\in U$, write $w=w_0v(\alpha)+w'$ with
$w'_0=0$, and note $D_\alpha w'=D_\alpha w$, so
$\norm{w'}\le(d-1)R^{d-2}\norm{D_\alpha w}$ as in the proof above. From
$0=\ell_\alpha^{\mathsf T}w=w_0P_1'(\alpha)+\ell_\alpha^{\mathsf T}w'$ we get
$\abs{w_0}\le\norm{\ell_\alpha}\norm{w'}/\sigma^{d-1}$, hence
$\norm w\le\theta_U\norm{D_\alpha w}$ with
$\theta_U=(d-1)R^{d-2}(1+dR^{d-1}d^2R^d\sigma^{-(d-1)})$. Therefore
$w^{\mathsf T}Q_*w\ge b_0^{d-2}\norm{D_\alpha w}^2\ge\beta_*\norm w^2$ on $U$,
where $\beta_*=b_0^{d-2}/\theta_U^2$.

\emph{The pencil.} Run the approximation with
$\delta\le\min\{1,\beta_*\}$ and let $B=Q$; then $B\in\mathcal L_{d-1}$ is
rational, $\norm{B-Q_*}\le\beta_*/2$, so $B|_U\succeq(\beta_*/2)I$ and
$\norm B\le\Lambda_*+1$. For $\lambda=(\lambda_1,\lambda_2)\in\R^2$ put
$Q_\lambda=C^{\mathsf T}BC-\lambda_1(C^{\mathsf T}B+BC)+\lambda_2B$ and
$q_\lambda(x)=(1,x)^{\mathsf T}Q_\lambda(1,x)$ for $x\in\R^{d-1}$. Since
$Cv(\alpha)=\alpha v(\alpha)$ and $v(\alpha)^{\mathsf T}Bv(\alpha)=0$, every
$q_\lambda$ vanishes at $p'$. At $\lambda_*=(\alpha,\alpha^2)$,
$Q_{\lambda_*}=(C-\alpha I)^{\mathsf T}B(C-\alpha I)$ is positive semidefinite
with kernel $\R v(\alpha)$, because $B$ is positive definite on
$U=\operatorname{image}(C-\alpha I)$. For $z$ with $z_0=0$, the first $d-1$
coordinates of $(C-\alpha I)z$ are $D_\alpha z$, so
$\norm{(C-\alpha I)z}\ge\norm z/((d-1)R^{d-2})$ and
\[
 z^{\mathsf T}Q_{\lambda_*}z\ge\gamma_G\norm z^2,\qquad
 \gamma_G=\frac{\beta_*}{2(d-1)^2R^{2d-4}} .
\]
Since $\norm C\le\norm C_F\le dR$, every $\lambda$ with
$\max\{\abs{\lambda_1-\alpha},\abs{\lambda_2-\alpha^2}\}\le3e$, where $e$ is a
dyadic number with $e\le\gamma_G/(6(2dR+1)(\Lambda_*+1))$, satisfies
$\norm{Q_\lambda-Q_{\lambda_*}}\le\gamma_G/2$, so $q_\lambda$ has Hessian
$2(Q_\lambda)_{1:d-1,1:d-1}\succeq\gamma_GI$.

\emph{A rational triangle.} Compute a dyadic $\hat\alpha$ with
$\abs{\hat\alpha-\alpha}\le e/(16(2R+1))$ (by Lemma~\ref{lem:algebraic-roots}
or by bisection), and put $c=(\hat\alpha,\hat\alpha^2)$; then
$\delta':=\lambda_*-c$ satisfies $\norm{\delta'}_\infty\le e/16$. Let
$\mu_1=c+(-e,-e)$, $\mu_2=c+(2e,-e)$, $\mu_3=c+(-e,2e)$. Solving
$\lambda_*=\sum_i\tau_i\mu_i$ with $\sum_i\tau_i=1$ gives
$\tau_2=(1+\delta'_1/e)/3$, $\tau_3=(1+\delta'_2/e)/3$, and
$\tau_1=1-\tau_2-\tau_3$, all larger than $1/4$. Every vertex lies within $3e$
of $\lambda_*$ in the maximum norm. Put $g_i=q_{\mu_i}$. These are rational
quadratics with positive definite Hessians that vanish at $p'$, and since
$\lambda\mapsto q_\lambda$ is affine,
\[
 \sum_i\tau_ig_i(x)=q_{\lambda_*}(x)=\omega^{\mathsf T}B\omega
 \ge\tfrac{\beta_*}2\norm\omega^2,\qquad \omega=(C-\alpha I)(1,x)\in U,
\]
with equality to zero only if $\omega=0$, that is, $(1,x)\in\R v(\alpha)$ and
$x=p'$. If
$g_i(x)\le0$ for all $i$, the left side is at most zero, so $x=p'$. Finally,
each $g_i$ is strictly convex with rational Hessian and gradient, so its
unique minimizer is rational and differs from the irrational point $p'$; its
minimum is therefore negative, and $\{g_i\le0\}$ is a full-dimensional
ellipsoid. All quantities have polynomial bit length, so the construction is
polynomial.
\end{proof}

\subsection{The cyclic family}
\label{app:algebraic-cyclic}

Throughout this subsection $n\ge2$, $m=n+1$, $\eta=(-1)^m$, indices are taken
modulo $m$, and $x_0=1$. Recall $d_n=(2^m-\eta)/3$, $e_i=((-2)^i-1)/3$, the
exponents $\theta_i$, the point $p_i=2^{e_i/d_n}$, and the binomials
\[
 r_i(x)=x_i^2-2^{\theta_i}x_{i+1}x_{i+2}\qquad(0\le i\le n).
\]
Write $a=2^{1/d_n}$, so $p_i=a^{e_i}$, and $w_i=p_i^{-2}$.

\begin{lemma}[Arithmetic of the cyclic point]
\label{lem:algebraic-cyclic-point}
\begin{enumerate}[label=(\alph*)]
\item $d_n$ is a positive odd integer, every $e_i$ is an integer, $e_0=0$,
  $e_1=-1$, and $\abs{e_i}<d_n$; hence $1/2<p_i<2$ and $1/4<w_i<4$.
\item $2e_i-e_{i+1}-e_{i+2}=d_n\theta_i$ for $0\le i\le n$, with indices of
  $e$ modulo $m$. Consequently $r_i(p)=0$ for all $i$.
\item $\Q(p)=\Q(p_1)=\Q(a)$ and $[\Q(a):\Q]=d_n$.
\item With $X_i=x_i/p_i$ (so $X_0=1$),
  $\sum_iw_ir_i(x)=\tfrac12\sum_i(X_i-X_{i+1})^2$. In particular $p$ is the
  only real common zero of the $r_i$.
\end{enumerate}
\end{lemma}

\begin{proof}
(a) $2^m\equiv(-1)^m\pmod3$, so $d_n$ is an integer; $2^m-\eta$ is odd, so
$d_n$ is odd. Likewise $(-2)^i\equiv1\pmod3$. For $i\le n$,
$\abs{e_i}\le(2^n+1)/3<(2^m-1)/3\le d_n$ because $n\ge2$.
(b) The sequence $E_i=((-2)^i-1)/3$, $i\ge0$, satisfies
$2E_i=E_{i+1}+E_{i+2}$, $E_m=\eta d_n$, and $E_{m+1}=-2\eta d_n-1$. For
$i\le m-3$ no index wraps and $\theta_i=0$. For $i=m-2$,
$2e_{m-2}-e_{m-1}-e_0=E_m=\eta d_n$. For $i=m-1$,
$2e_{m-1}-e_0-e_1=2E_{m-1}+1=E_m+E_{m+1}+1=-\eta d_n$. Since $a^{d_n}=2$,
$p_i^2=a^{d_n\theta_i}p_{i+1}p_{i+2}=2^{\theta_i}p_{i+1}p_{i+2}$.
(c) Every $p_i$ is an integer power of $a$, and $p_1=a^{-1}$. The polynomial
$T^{d_n}-2$ is irreducible by Eisenstein's criterion at $2$.
(d) By (b), $w_ir_i=X_i^2-X_{i+1}X_{i+2}$; summing over $i$ modulo $m$ gives
$\sum_iX_i^2-\sum_iX_iX_{i+1}=\frac12\sum_i(X_i-X_{i+1})^2$. This vanishes
only if all $X_i$ equal $X_0=1$.
\end{proof}

\begin{lemma}[Quantitative data]
\label{lem:algebraic-cyclic-data}
Write $r_i(p+u)=b_i^{\mathsf T}u+u^{\mathsf T}T_iu$ and let $J$ have rows
$b_i^{\mathsf T}$. Then $\norm{T_i}\le1$, $\norm{b_i}<7$, $\norm J<32$,
$J^{\mathsf T}J\succeq(8n)^{-2}I$, and
$G_*=\sum_iw_ir_i$ satisfies $G_*(p+u)=u^{\mathsf T}H_*u$ with
$(8n^2)^{-1}I\preceq H_*\preceq8I$.
\end{lemma}

\begin{proof}
Because $m\ge3$, the indices $i,i+1,i+2$ are distinct, so the quadratic part
of $r_i$ is a diagonal entry $1$ (absent if $i=0$) and a disjoint off-diagonal
pair $-2^{\theta_i}/2$ (absent if an index is $0$); its norm is at most $1$.
The gradient $b_i$ has at most three nonzero entries, $2p_i$,
$-2^{\theta_i}p_{i+2}$, and $-2^{\theta_i}p_{i+1}$, each of absolute value less
than $4$. Put $h=ED_p^{-1}u$, where $D_p=\diag(p_1,\ldots,p_n)$ and $E$ inserts
a zero coordinate with index $0$. Using $2^{\theta_i}p_{i+1}p_{i+2}=p_i^2$,
\[
 b_i^{\mathsf T}u=p_i^2(2h_i-h_{i+1}-h_{i+2}),\qquad
 J=\diag(p_0^2,\ldots,p_n^2)\,(2I-\mathcal S-\mathcal S^2)\,ED_p^{-1},
\]
where $(\mathcal Sz)_i=z_{i+1}$ is the cyclic shift. The four factors have
norms less than $4$, at most $4$, $1$, and less than $2$, so $\norm J<32$.
For the lower bound, $2I-\mathcal S-\mathcal S^2=(2I+\mathcal S)(I-\mathcal S)$
and $\norm{(2I+\mathcal S)y}\ge2\norm y-\norm{\mathcal Sy}=\norm y$. Since
$h_0=0$, each $h_k$ is a sum of $k\le n$ consecutive differences, so by the
Cauchy--Schwarz inequality
\begin{equation}
 \norm h^2\le n^2\sum_i(h_i-h_{i+1})^2=n^2\norm{(I-\mathcal S)h}^2 .
 \label{eq:algebraic-anchored}
\end{equation}
With $\norm h\ge\norm u/2$ and $p_i^2>1/4$ this gives
$\norm{Ju}\ge\norm u/(8n)$. Finally, Lemma~\ref{lem:algebraic-cyclic-point}(d)
with $X_i=1+h_i$ gives $G_*(p+u)=\tfrac12\norm{(I-\mathcal S)h}^2$, which lies
between $\norm h^2/(2n^2)\ge\norm u^2/(8n^2)$ and $2\norm h^2\le8\norm u^2$.
\end{proof}

\begin{lemma}[Rational rank-one square roots]
\label{lem:algebraic-rank-one}
Let $\mathbf g\in\Q^m$ be nonzero, $R_{\mathbf g}=\mathbf g^{\mathsf T}\mathbf g$,
and let $0<\varrho<\sqrt{R_{\mathbf g}}$ be rational. Put
$t=(R_{\mathbf g}-\varrho^2)/(2\varrho)$ and
$L_\varrho=tI+(\varrho/R_{\mathbf g})\mathbf g\mathbf g^{\mathsf T}$. Then
$t>0$, $L_\varrho$ is rational, symmetric and invertible, and
$L_\varrho^2=t^2I+\mathbf g\mathbf g^{\mathsf T}$. Hence for every vector
$r=(r_0,\ldots,r_{m-1})$ of polynomials,
$(\mathbf g^{\mathsf T}r)^2+t^2\sum_ir_i^2=\sum_i(L_\varrho r)_i^2$.
\end{lemma}

\begin{proof}
Since $(\mathbf g\mathbf g^{\mathsf T})^2=R_{\mathbf g}\mathbf g\mathbf
g^{\mathsf T}$, the coefficient of $\mathbf g\mathbf g^{\mathsf T}$ in
$L_\varrho^2$ is $(2t\varrho+\varrho^2)/R_{\mathbf g}=1$. The eigenvalues of
$L_\varrho$ are $t$ and $t+\varrho$.
\end{proof}

\begin{proof}[Proof of Theorem~\ref{thm:algebraic-cyclic}]
\emph{Integer weights.} Let $M_1=10^6n^5$ and $Q_1=32mM_1^2$. We compute
integers $z_0,\ldots,z_n$ with $\abs{z_i/Q_1-w_i}\le1/Q_1$ in time polynomial
in $n$, without ever forming a polynomial of degree $d_n$. Since
$w_i=\exp(y_i)$ with $y_i=-(2e_i/d_n)\log2\in(-2\log2,2\log2)$, use
$\log2=2\sum_{j\ge0}3^{-(2j+1)}/(2j+1)$, whose tail after $N$ terms is less
than $9^{-N}$. Replacing $\log2$ by the $N$-term sum $L_N$ changes $y_i$ by
less than $2\cdot9^{-N}$, and $\exp$ is $8$-Lipschitz on $[-2,2]$. The
Taylor polynomial of degree $K\ge2$ of $\exp$ approximates it on $[-2,2]$
within $2^{K+2}/(K+1)!$. Choosing $N$ and $K$ of size $O(\log Q_1)$ gives a
rational $\hat w_i$ with $\abs{\hat w_i-w_i}\le1/(4Q_1)$; rounding $Q_1\hat
w_i$ to an integer gives $z_i$. The integers $e_i$ and $d_n$ have $O(n)$ bits,
and all intermediate rationals have bit length polynomial in $n$.

\emph{The exposing quadratic.} Let $\mathbf g=(z_i/Q_1)_{i=0}^n$ and
$G=\mathbf g^{\mathsf T}r=\sum_i(z_i/Q_1)r_i$. Then $G(p)=0$, and with
$G(p+u)=\ell^{\mathsf T}u+u^{\mathsf T}Hu$, the identity
$\sum_iw_ib_i=\nabla G_*(p)=0$ gives $\ell=\sum_i(z_i/Q_1-w_i)b_i$ and
$H-H_*=\sum_i(z_i/Q_1-w_i)T_i$. By Lemma~\ref{lem:algebraic-cyclic-data},
\[
 \norm\ell\le\frac{8m}{Q_1}=\frac1{4M_1^2},\qquad
 \norm{H-H_*}\le\frac m{Q_1}=\frac1{32M_1^2},\qquad
 \frac1{16n^2}I\preceq H\preceq9I .
\]

\emph{Regularization with a rational square root.} Since $1/5<z_i/Q_1<5$, we
have $m/25<R_{\mathbf g}<25m$. The function
$\varrho\mapsto t(\varrho)=(R_{\mathbf g}-\varrho^2)/(2\varrho)$ decreases on
$(0,\infty)$, and the set of $\varrho$ with $1/M_1\le t(\varrho)\le2/M_1$ is an
interval of length at least
$M_1^{-1}\min_{1/M_1\le t\le2/M_1}R_{\mathbf g}/\bigl(\sqrt{R_{\mathbf
g}+t^2}(\sqrt{R_{\mathbf g}+t^2}+t)\bigr)\ge1/(1450M_1)$ contained in
$(0,5m)$. Rational bisection on $[0,5m]$, testing $M_1(R_{\mathbf
g}-\varrho^2)\ge2\varrho$ and $M_1(R_{\mathbf g}-\varrho^2)\le4\varrho$ at
dyadic midpoints, finds such a $\varrho=a_1/b_1$ after $O(\log n)$ steps, with
$a_1,b_1$ of polynomial magnitude. Put $t=t(\varrho)$ and $\varepsilon=t^2$,
so $M_1^{-2}\le\varepsilon\le4M_1^{-2}$.

\emph{Curvature and Gram.} Apply Lemma~\ref{lem:quartic-realization} to $G$ and
the $m=n+1$ residuals $r_0,\ldots,r_n$ with $\mu=1/(16n^2)$, $\Lambda=9$,
$\beta=8$, and $\nu=1/(8n)$. Here $\Lambda+m\beta=8n+17\le17n$, and
\[
 \frac{\mu^2}{2m}\ge\frac1{768n^5},\qquad
 \frac{\nu^2\mu^2}{36n(\Lambda+m\beta)^2}\ge\frac1{170459136\,n^9},
\]
both larger than $4M_1^{-2}=4\cdot10^{-12}n^{-10}$; also
$\norm\ell\le\varepsilon/4$. Thus $\Phi=G^2+\varepsilon\sum_ir_i^2$ satisfies
$\nabla^2\Phi\succeq\frac32\varepsilon\nu^2I$, and the canonical Gram of this
representation is positive definite. By Lemma~\ref{lem:algebraic-rank-one},
$\Phi=\sum_i(L_\varrho r)_i^2$. Both representations equal
$r^{\mathsf T}(\mathbf g\mathbf g^{\mathsf T}+\varepsilon I)r$ with the same
matrix, so they have the same canonical Gram.

\emph{Integer squares.} Let $S_0=\sum_iz_i^2$ and $D_0=2a_1b_1Q_1^2S_0$. Then
$L_\varrho=tI+\frac{a_1}{b_1S_0}zz^{\mathsf T}$ with
$t=(S_0b_1^2-a_1^2Q_1^2)/(2a_1b_1Q_1^2)$, so $D_0L_\varrho$ is an integer
matrix; each $2r_i$ has integer coefficients. Put
$s_i=32M_1nD_0(L_\varrho r)_i=16M_1n\sum_j(D_0L_\varrho)_{ij}(2r_j)\in\Z[x]$
and $F_n=\sum_is_i^2=(32M_1nD_0)^2\Phi$. The integers used in this
denominator-clearing and scaling step, namely $z_i$, $Q_1$, $M_1$, $a_1$,
$b_1$, $S_0$, $D_0$, and the coefficients of the $2r_j$, have magnitude
polynomial in $n$, so the coefficients of the $s_i$ have $O(\log(n+1))$
bits. The integers $d_n$ and $e_i$, which have $O(n)$ bits, enter only
through the weights $w_i$ and are never coefficients. Since $D_0\ge1$ and
$\varepsilon\ge M_1^{-2}$,
$\nabla^2F_n\succeq(32M_1n)^2\cdot\frac32M_1^{-2}(8n)^{-2}I=24I$.

The canonical Gram of $\sum_is_i^2$ is $(32M_1nD_0)^2$ times that of $\Phi$,
hence positive definite. It is an integer matrix: for an integer quadratic
$s=s_0+b^{\mathsf T}x+x^{\mathsf T}Tx$ the matrices $b$, $2T$ are integral, and
every entry of \eqref{eq:algebraic-canonical} is an integer combination of
products of entries of $b$, $2T$ and $s_0$. This proves (a). Part (b) is
Lemma~\ref{lem:algebraic-cyclic-point}, since $F_n$ vanishes exactly at the
common zeros of the $r_i$ ($L_\varrho$ is invertible).

(c) $G$ has $O(n)$ monomials, so $G^2$ and $\sum_ir_i^2$ have $O(n^2)$, and
every coefficient of $F_n=\sum_is_i^2$ is a sum of $O(n^2)$ products of
polynomial-magnitude integers.

(d) is Theorem~\ref{thm:sos-length}, which applies because $F_n$ is globally
convex of degree four with $\nabla^2F_n(p)\succ0$.

(e) Since all factors vanish at $p$,
$\nabla^2\Phi(p)=2\ell\ell^{\mathsf T}+2\varepsilon J^{\mathsf T}J$. Its least
eigenvalue is at least $2\varepsilon\nu^2=\varepsilon/(32n^2)$ and its
largest at most $2\norm\ell^2+2\varepsilon\norm J^2\le\varepsilon(\varepsilon/8
+2048)$. The ratio is at most $32n^2(1/8+2048)<65600n^2$, and scaling by a
positive constant does not change it.

(f) Substituting $x_1=y_1-1$ maps each $s_i$ to an integer quadratic with
$O(\log(n+1))$-bit coefficients; the curvature bound, the integer canonical
Gram (Lemma~\ref{lem:algebraic-gram}(b)), and the number of squares are
unchanged. The zero has first coordinate $1+a^{-1}$. The number $a^{-1}$ is a
root of $2T^{d_n}-1$, which is irreducible because its reversal $T^{d_n}-2$ is;
so $1+a^{-1}$ has the irreducible primitive minimal polynomial
$2(T-1)^{d_n}-1$, primitive because its leading coefficient is $2$ and its
constant coefficient is $2(-1)^{d_n}-1=-3$. The coefficient of $T^j$ for
$j\ge1$ is $\pm2\binom{d_n}j\ne0$. Each coefficient has $O(d_n)$ bits, and for
the $\Theta(d_n)$ indices $d_n/3\le j\le2d_n/3$ one has
$\binom{d_n}j\ge2^{\min(j,d_n-j)}\ge2^{d_n/3-1}$; hence the total binary length
is $\Theta(d_n^2)$.
\end{proof}

\begin{remark}[Rational ellipsoids for the cyclic point]
\label{rem:algebraic-cyclic-ellipsoids}
With $\tau_e=1/(32n^2)$, $c_i=z_i/Q_1$, $C_z=\sum_ic_i$ and $W=\sum_iw_i$, the
$n+1$ rational quadratics $G+\tau_er_i$ have quadratic parts at least
$(\mu-\tau_e)I=I/(32n^2)$ and vanish at $p$. The real weights
$\lambda_i=\tau_e^{-1}(w_i-Wc_i/(\tau_e+C_z))$ satisfy
$\sum_i\lambda_i(G+\tau_er_i)=G_*$, and
$\tau_e(\tau_e+C_z)\lambda_i=\tau_ew_i+w_i\sum_j(c_j-w_j)-W(c_i-w_i)\ge\tau_e/4-8m/Q_1>0$.
Since $G_*\ge0$ vanishes only at $p$, the intersection of the ellipsoids
$\{G+\tau_er_i\le0\}$ is $\{p\}$. Each ellipsoid is full-dimensional because
its rational minimizer differs from $p$.
\end{remark}

\subsection{The minimum number of squares}
\label{app:algebraic-length}

We use the Brouwer degree of continuous maps $S^r\to S^r$, $r\ge1$: it is a
homotopy invariant, and if a continuous map is smooth near the finite
preimage of a point $y$ at which it is a local diffeomorphism, its degree is
the sum over that preimage of the signs of the Jacobian determinants
\cite[Proposition~2.30]{Hatcher2002}, \cite[Sections~4--5]{Milnor1965}. A
continuous map $q:\R^r\to\R^r$ with $\norm{q(w)}\to\infty$ as
$\norm w\to\infty$ extends continuously to the one-point compactification
$S^r$ by $\infty\mapsto\infty$; its degree is that of the extension. A
homotopy $q_s$, $0\le s\le1$, extends to a homotopy of the extensions when
$\norm{q_s(w)}\ge\psi(\norm w)$ for a function $\psi$ with
$\psi(\rho)\to\infty$, uniformly in $s$. Sard's theorem gives regular values
in every nonempty open set \cite[Section~3]{Milnor1965}.

\begin{proof}[Proof of Theorem~\ref{thm:sos-length}]
\emph{Step 1: degrees and the count $m\ge n$.} If some $g_i$ had degree
$e>2$, the sum of the squares of the degree-$e$ parts of the factors of
maximal degree would be a nonzero form of degree $2e>4$ (a nonzero real form
squared is nonzero, and squares cannot cancel), contradicting $\deg f=4$. So
every $g_i$ is a quadratic. Write $q=(g_1,\ldots,g_m)$. All $g_i$ vanish at
$p$, so $\nabla^2f(p)=2Dq(p)^{\mathsf T}Dq(p)$; positive definiteness forces
$\rank Dq(p)=n$ and $m\ge n$. The zero $p$ is unique: if $f(p')=0$ with
$p'\ne p$, then by convexity and $f\ge0$, $f$ vanishes on the segment
$[p,p']$, so the second derivative of $f$ at $p$ in the direction $p'-p$ is
zero. Suppose now, for a contradiction, that $m=n$, so $q:\R^n\to\R^n$.

\emph{Step 2: the quartic-flat subspace.} Write $q=q_2+q_1+q_0$ by
homogeneous degree and $H(x)=\norm{q_2(x)}^2$, the quartic form of $f$. The
function $H$ is convex, as the pointwise limit of $s^{-4}f(sx)$ when
$s\to\infty$, and it is nonnegative and even. Hence
$K=\{x:H(x)=0\}$ is convex and closed under real scalar multiples, so it is a
linear subspace. For $v\in K$, $x\in\R^n$, $s\in\R$, and $0<\eta<1$,
\[
 H(x+sv)\le(1-\eta)H\Bigl(\frac x{1-\eta}\Bigr)+\eta H\Bigl(\frac{sv}\eta\Bigr)
 =(1-\eta)^{-3}H(x).
\]
Letting $\eta\to0$ gives $H(x+sv)\le H(x)$, and applying this to $x+sv$ and
$-s$ gives equality. Let $B_2$ be the symmetric bilinear map with
$q_2(x)=B_2(x,x)$. Since $q_2(v)=0$, $q_2(x+sv)=q_2(x)+2sB_2(x,v)$, and the
coefficient of $s^2$ in the constant function $\norm{q_2(x+sv)}^2$ is
$4\norm{B_2(x,v)}^2$. Hence $B_2(x,v)=0$, and $q_2(x+v)=q_2(x)$ for all
$v\in K$.

\emph{Step 3: coordinates.} Choose orthonormal coordinates $x=(w,v)$ with
$v\in K\cong\R^k$ and $w\in K^\perp\cong\R^{r}$, $r=n-k$. Then
\[
 q(w,v)=q_2(w)+Aw+Bv+c
\]
with matrices $A$, $B$ and a vector $c$. The matrix $Dq(p)=(Dq_2(p_w)+A\ \ B)$
is nonsingular, so $B$ has full column rank $k$. Because $\deg f=4$,
$q_2\ne0$ and $r\ge1$. For $w\ne0$ in $K^\perp$, $w\notin K$, so
$\norm{q_2(w)}>0$.

\emph{Step 4: convexity removes the coupling.} The $ww$ block of the Hessian
of $f=\sum_iq_i^2$ is
\[
 \nabla^2_{ww}f(w,v)=C(w)+2\sum_i(Bv)_i\nabla^2q_{2,i},
\]
where $C(w)=\sum_i\bigl[2\nabla_wq_i\nabla_wq_i^{\mathsf
T}+2(q_{2,i}(w)+(Aw)_i+c_i)\nabla^2q_{2,i}\bigr]$ does not depend on $v$,
because $\nabla_wq_i=\nabla q_{2,i}(w)+A_i^{\mathsf T}$. Global convexity
makes this block positive semidefinite for all $v$. Taking $v=se_j$ with
$s\in\R$ arbitrary shows that the symmetric matrix
$\sum_i(Be_j)_i\nabla^2q_{2,i}$ vanishes, since otherwise
$y^{\mathsf T}(C(w)+2s\sum_i(Be_j)_i\nabla^2q_{2,i})y<0$ for a suitable $y$
and $s$. So the homogeneous quadratic $(Be_j)^{\mathsf T}q_2(w)$ has zero
Hessian and is identically zero:
\begin{equation}
 B^{\mathsf T}q_2(w)=0\qquad(w\in\R^r).
 \label{eq:algebraic-length-coupling}
\end{equation}

\emph{Step 5: elimination of $v$.} Let $T_\perp$ be an $r\times n$ matrix whose
rows form an orthonormal basis of $(\operatorname{im}B)^\perp$, and put
\[
 v_0(w)=-(B^{\mathsf T}B)^{-1}B^{\mathsf T}(Aw+c),\qquad
 \bar q(w)=T_\perp\bigl(q_2(w)+Aw+c\bigr).
\]
If $k=0$, take $T_\perp=I$ and omit $v$. By
\eqref{eq:algebraic-length-coupling}, $q_2(w)\in(\operatorname{im}B)^\perp$, so
orthogonal decomposition of $q(w,v)$ along $\operatorname{im}B$ and its
complement gives
\[
 f(w,v)=\norm{B(v-v_0(w))}^2+\norm{\bar q(w)}^2 .
\]
The function $\bar f(w)=f(w,v_0(w))=\norm{\bar q(w)}^2$ is convex, being the
restriction of $f$ to the graph of the affine map $v_0$. Its zeros lift to
zeros of $f$, so its only zero is $w_*=p_w$, and $p_v=v_0(p_w)$. With
$E=\left(\begin{smallmatrix}I_r\\Dv_0\end{smallmatrix}\right)$ of full column
rank, $\nabla^2\bar f(w_*)=E^{\mathsf T}\nabla^2f(p)E\succ0$; as
$\nabla^2\bar f(w_*)=2D\bar q(w_*)^{\mathsf T}D\bar q(w_*)$, the square map
$\bar q:\R^r\to\R^r$ has the unique zero $w_*$, and it is regular. Its
quadratic part $\bar q_2=T_\perp q_2$ satisfies
$\norm{\bar q_2(w)}=\norm{q_2(w)}$, which is positive on the unit sphere; by
compactness and homogeneity there is $a_0>0$ with
$\norm{\bar q_2(w)}\ge a_0\norm w^2$.

\emph{Step 6: parity.} For $0\le s\le1$ put
$\bar q^{(s)}(w)=\bar q_2(w)+s(T_\perp Aw+T_\perp c)$. Then
$\norm{\bar q^{(s)}(w)}\ge a_0\norm w^2-\norm{T_\perp A}\norm w-\norm{T_\perp c}$
uniformly in $s$, so $\bar q=\bar q^{(1)}$ and $\bar q_2=\bar q^{(0)}$ have
the same degree. The point $0$ has the single preimage $w_*$ under $\bar q$,
where $\bar q$ is a local diffeomorphism, so $\deg\bar q=\pm1$. On the other
hand, choose a regular value $y\ne0$ of $\bar q_2$. Its preimage is compact,
by the growth bound, and discrete, so it is finite; it does not contain $0$,
and it is invariant under $w\mapsto-w$ because $\bar q_2$ is even. Since
$D\bar q_2$ is linear in $w$, $D\bar q_2(-w)=-D\bar q_2(w)$, and the two
points of each pair $\{w,-w\}$ contribute $\epsilon$ and $(-1)^r\epsilon$ for
some $\epsilon=\pm1$. Each pair contributes an even number, so
$\deg\bar q_2$ is even. This contradiction proves $m\ge n+1$.

\emph{Sharpness.} For $f=(\sum_ix_i^2)^2+\sum_ix_i^2$,
$\nabla^2f=8xx^{\mathsf T}+(4\norm x^2+2)I\succeq2I$, $f(0)=0$, and $f$ is a sum
of $n+1$ squares. Lemma~\ref{lem:quartic-realization} with $m=n$ and
Theorem~\ref{thm:algebraic-cyclic} give further examples with exactly $n+1$
squares. The examples after Theorem~\ref{thm:sos-length} are verified
directly; for the sextic,
$\frac{d^2}{dx_1^2}(x_1+x_1^3)^2=2+24x_1^2+30x_1^4$.
\end{proof}

\subsection{Binomial networks}
\label{app:algebraic-network}

We use the following classical facts about a directed multigraph
$\mathcal D$ on $\{0,\ldots,n\}$ with out-degree two everywhere, $A$ its
adjacency matrix (counting multiplicities and loops), and $L=2I-A$. Parallel
arcs are distinguished.
\begin{enumerate}[label=(N\arabic*)]
\item (Directed matrix-tree theorem \cite[Theorem~3.6]{Tutte1948}.) The principal minor
  of $L$ obtained by deleting row and column $i$ equals the number $\tau_i$ of
  spanning arborescences in which every vertex has a directed path to $i$.
  Loops cannot enter a tree, and their contributions to the degree and
  adjacency terms of $L$ cancel; hence the theorem includes loops here.
\item (BEST theorem \cite{vanAardenneEhrenfestDeBruijn1951}.) If $\mathcal D$
  is Eulerian (every in-degree is also two), the number of Euler circuits,
  counted as cyclic arc sequences, equals
  $\tau_i\prod_v(\deg^+v-1)!=\tau_i$ for every $i$.
\item (Interlace polynomials \cite{ArratiaBollobasSorkin2004}.) If $\mathcal D$
  is Eulerian with Euler circuit $\mathcal C$, the interlace graph $H$ of
  $\mathcal C$ on the $m=n+1$ vertices has an interlace polynomial
  $q_H(x)=\sum_{j\ge1}\alpha_jx^j$ with nonnegative integer coefficients and no
  constant term, $q_H(1)$ equals the number of Euler circuits of $\mathcal D$,
  and $q_H(2)=2^m$.
\item (\cite[Theorem~1]{BalisterBollobasCutlerPebody2002}.) For a graph $H$ on
  $m$ vertices with adjacency matrix $A_H$,
  $q_H(-1)=(-1)^{\mathrm{rk}}2^{m-\mathrm{rk}}$, where
  $\mathrm{rk}=\rank_{\mathbb F_2}(I+A_H)$.
\end{enumerate}

\begin{proof}[Proof of Proposition~\ref{prop:algebraic-network}]
\emph{Lattice count.} Let $B$ be $L$ with column $0$ deleted; its rows are the
exponent vectors of the Laurent monomials $y_i^2/(y_{a(i)}y_{b(i)})$ in
$y_1,\ldots,y_n$ (with $y_0=1$). Strong connectivity gives $\tau_i>0$, so by
(N1) $\rank L=m-1$. As $L\mathbf1=0$ and
$L\operatorname{adj}(L)=\operatorname{adj}(L)L=\det(L)I=0$, the columns of
$\operatorname{adj}(L)$ lie in $\R\mathbf1$ and its rows in the
one-dimensional left kernel; its diagonal is $(\tau_i)$, so
$\operatorname{adj}(L)=\mathbf1\tau^{\mathsf T}$ and $\tau^{\mathsf T}L=0$. The
maximal minor of $B$ that omits row $i$ is $\pm\operatorname{adj}(L)_{0i}=
\pm\tau_i$. By the Smith normal form, the index of the row lattice
$\Lambda_B$ of $B$ in $\Z^n$ equals the gcd of the maximal minors,
$g=\gcd(\tau_0,\ldots,\tau_n)$.

\emph{Solutions.} Every complex solution has nonzero coordinates: $x_0=1$,
and if $x_i\ne0$ then $c_ix_{a(i)}x_{b(i)}=x_i^2\ne0$; strong connectivity
propagates this to all vertices. Writing $x_i=p_iy_i$, the equations become
$y^{B_i}=1$ for every row $B_i$, because $p^{B_i}=c_i$. Thus normalized
solutions correspond to characters of the finite abelian group
$\Gamma_B=\Z^n/\Lambda_B$, of which there are exactly $g$. If $\epsilon$ is the
exponent of $\Gamma_B$, then $\epsilon e_j=\sum_ik_iB_i$ with integers $k_i$,
so $p_j^\epsilon=\prod_ic_i^{k_i}\in\Q$; hence $p$ is algebraic, every
embedding of $\Q(p)$ gives a distinct complex solution, and
$[\Q(p):\Q]\le g$. Real normalized solutions are characters with values in
$\{\pm1\}$, the only real roots of unity; there is a nontrivial one exactly
when $g$ is even. Since $p$ is the only real solution, $g$ is odd.

\emph{Non-Eulerian graphs.} The primitive positive integer vector
$\omega=\tau/g$ spans the left kernel of $L$. It equals $\mathbf1$ exactly when
$\mathbf1^{\mathsf T}L=0$, that is, when every in-degree is two. Otherwise some
$\omega_i\ge2$. An arborescence toward $i$ chooses one of the two outgoing arcs
at each of the $n$ other vertices, so $\tau_i\le2^n$. Hence
$g\le\tau_i/2\le2^{n-1}$, and since $g$ is odd, $g\le2^{n-1}-1<d_n$.

\emph{Eulerian graphs.} Now all $\tau_i$ equal $g$, which by (N2) and (N3)
equals $q_H(1)$ for the interlace graph of any Euler circuit. Since
$q_H(1)-q_H(-1)=2\sum_{j\text{ odd}}\alpha_j$ is even and $q_H(1)=g$ is odd,
$q_H(-1)$ is odd, so (N4) forces $\mathrm{rk}=m$ and $q_H(-1)=(-1)^m$. For
$j\ge1$, $2^j-(-1)^j\ge3$, so
\[
 3g\le\sum_j\alpha_j\bigl(2^j-(-1)^j\bigr)=q_H(2)-q_H(-1)=2^m-(-1)^m,
\]
that is, $g\le d_n$.

\emph{The exposing normal form.} Let $\mathcal Q=\sum_j\lambda_jr_j$ be the
positive semidefinite combination with kernel $\R p$, $p=(1,p_1,\ldots,p_n)$.
Every diagonal entry of $\mathcal Q$ is positive: a zero diagonal entry would
put a coordinate vector into the kernel, but $p$ has no zero coordinate and
$m\ge2$. A nonzero binomial vanishing at $p$ has two distinct monomials with
nonzero coefficients (a single monomial cannot vanish at $p$), and it cannot
have two square monomials with positive coefficients. So each summand
$\lambda_jr_j$ contributes a positive coefficient to at most one square
$X_i^2$, while all $m$ diagonal entries must receive one. Hence each summand
supplies exactly one positive square, at distinct indices. Assigning to the
binomial its square $X_i^2$ and rescaling rationally gives
$X_i^2-c_iX_{a(i)}X_{b(i)}$ with $c_i\in\Q\setminus\{0\}$, which defines a
two-out multigraph. The same count shows that fewer than $m$ binomials cannot
supply all diagonal entries.

In the coordinates $X_i=p_iY_i$, the $i$-th summand becomes
$h_i(Y_i^2-Y_{a(i)}Y_{b(i)})$ with $h_i>0$. The kernel contains
$\mathbf1$, which gives $2h=A^{\mathsf T}h$: every vertex has inflow equal to
outflow when arc $i\to j$ carries flow $h_i$. If $\mathcal D$ were not
strongly connected, take a strongly connected component with no incoming
arcs from other components; summing the flow balance over it shows that its
outgoing arcs carry zero flow, so it has none, and the argument repeats on
the rest. Thus the components are mutually disconnected. By stationarity,
\[
 \sum_ih_i\bigl(Y_i^2-Y_{a(i)}Y_{b(i)}\bigr)
 =\tfrac12\sum_ih_i\bigl(Y_{a(i)}-Y_{b(i)}\bigr)^2,
\]
because the diagonal coefficients agree, so the indicator vector of each
component lies in the kernel. Corank one forces a single component. Finally,
the common real zero is unique: if all binomials vanish at $(1,x)$, then
$\mathcal Q(1,x)=0$ and $(1,x)\in\R p$, so $x=(p_1,\ldots,p_n)$. The first part
now applies.
\end{proof}

\subsection{Imported results from algebraic geometry}
\label{app:algebraic-imports}

All varieties are over $\C$ and projective unless stated otherwise; a variety
is reduced and irreducible, and its degree is the number of points in its
intersection with a general linear space of complementary dimension
\cite{Harris1992}. $S=\C[x_0,\ldots,x_N]$, and for a
zero-dimensional subscheme $R\subset\mathbb P^N$ and $k\ge0$,
$H_R(k)$ is the rank of the restriction map $S_k\to H^0(\mathcal O_R(k))$.
\begin{enumerate}[label=(AG\arabic*)]
\item \cite[Propositions~I.7.1, I.7.2 and Theorem~I.7.7]{Hartshorne1977} If
  $X\subseteq\mathbb P^N$ is a variety of dimension $r\ge1$ and $V$ is a
  hypersurface of degree $e$ not containing $X$, then $X\cap V$ is nonempty,
  every irreducible component of $X\cap V$ has dimension $r-1$, and the sum
  of their degrees is at most $e\deg X$.
\item \cite[Proposition~8.4 and Example~8.4.6]{Fulton1998} If forms
  $F_1,\ldots,F_N$ of degrees $e_1,\ldots,e_N$ have finitely many common zeros
  in $\mathbb P^N$, then they form a regular sequence, the scheme
  $Z=V(F_1,\ldots,F_N)$ (a \emph{complete intersection}) has length
  $e_1\cdots e_N$, $S/(F_1,\ldots,F_N)$ is Cohen--Macaulay of dimension one, and
  every linear form that vanishes at no point of $Z$ is a nonzerodivisor on
  $S/(F_1,\ldots,F_N)$.
\item (Cayley--Bacharach \cite[Theorem~CB7]{EisenbudGreenHarris1996}.)
  Let $N\ge3$, let $Z\subset\mathbb P^N$ be a zero-dimensional complete
  intersection of hypersurfaces of degrees $e_1,\ldots,e_N$, and put
  $s=\sum_ie_i-N-1$. If $Z$ is reduced, every form of degree $s$ vanishing at
  all but one point of $Z$ vanishes on $Z$. In general, if $\Gamma,R\subseteq Z$
  are closed subschemes residual to each other, then for every integer $t$
  with $0\le t\le s$,
  $h^0(\mathcal I_\Gamma(t))-h^0(\mathcal I_Z(t))=h^1(\mathcal I_R(s-t))$.
  Moreover, for every finite subscheme $R\subset\mathbb P^N$ and every
  $k\ge0$, $h^1(\mathcal I_R(k))=\operatorname{length}(R)-H_R(k)$, by the
  cohomology sequence of
  $0\to\mathcal I_R(k)\to\mathcal O(k)\to\mathcal O_R(k)\to0$ and
  $H^1(\mathbb P^N,\mathcal O(k))=0$. We use the reduced statement for
  $N=n\ge3$ and the residual statement only for $N=n\ge4$, with $t=2$ and
  $s-t=n-3\ge1$.
\item (Gorenstein duality \cite[Chapter~21]{Eisenbud1995}.) If
  $\bar q_1,\ldots,\bar q_N\in\C[y_1,\ldots,y_N]$ is a regular sequence of
  quadratic forms, the graded algebra $\mathcal G=\C[y]/(\bar q)$ has Hilbert
  series $(1+t)^N$, $\mathcal G_N$ is one-dimensional, and multiplication
  $\mathcal G_i\times\mathcal G_{N-i}\to\mathcal G_N$ is a perfect pairing. The
  same holds over $\Q$.
\item (Residual intersections
  \cite[Theorem~3.2 and Remark~3.3]{EklundJostPeterson2013}; Segre classes of
  local complete intersections \cite[Proposition~4.1]{Fulton1998}.) Let
  $Y\subset\mathbb P^5$ be a nonempty reduced pure-dimensional local complete
  intersection, not necessarily smooth (for example a reduced curve that is a
  complete intersection of type $(1,1,2,2)$ with singular points), such that
  $\mathcal I_Y(2)$ is generated by global sections. Then for general quadrics
  $h_1,\ldots,h_5$ vanishing on $Y$, the set $V(h_1,\ldots,h_5)\setminus Y$ is
  finite, and the sum of the local lengths of $V(h_1,\ldots,h_5)$ at the
  points of $V(h_1,\ldots,h_5)\setminus Y$ equals
  $32-\int_Y(1+2H)^5\cap s(Y,\mathbb P^5)$. For a local complete
  intersection, singular or not, the Segre class is $s(Y,\mathbb
  P^5)=c(N_{Y/\mathbb P^5})^{-1}\cap[Y]$, with $N_{Y/\mathbb P^5}$ the normal
  bundle, so the count is $32-e(Y)$ with
  $e(Y)=\int_Y\bigl\{(1+2H)^5c(N_{Y/\mathbb P^5})^{-1}\bigr\}_{\dim Y}$. The
  cited residual count holds for an arbitrary closed subscheme cut out by
  forms of degree at most two; the local complete intersection property is
  used only for this formula for the Segre class.
\item (Minimal degree \cite[Proposition~0, Theorem~1 and
  Lemma~2.1]{EisenbudHarris1987}.) An irreducible variety of dimension $k$ that
  spans $\mathbb P^r$ has degree at least $r-k+1$. An irreducible curve of
  degree $r$ spanning $\mathbb P^r$ is a rational normal curve, which is smooth
  of genus zero and whose ideal is generated by quadrics.
\item A smooth curve $C\subset\mathbb P^5$ of degree $\delta$ and genus
  $\mathrm g$ has $\deg N_{C/\mathbb P^5}=6\delta+2\mathrm g-2$, by the Euler
  sequence and adjunction; consequently $e(C)=4\delta-2\mathrm g+2$.
\end{enumerate}

We also use two elementary consequences. First, a variety $X$ that is
invariant under complex conjugation and has odd degree has a real point:
intersecting with general real linear subspaces of complementary dimension
gives $\deg X$ distinct points, permuted by conjugation, and nonreal points
pair up. In particular a conjugation-invariant linear space has real points.
Second, a finite subscheme of $\mathbb P^N$ defined over $\R$ and of odd length
has a real point, because conjugation pairs nonreal points with equal local
lengths.

\begin{lemma}[Weighted B\'ezout count]
\label{lem:algebraic-weighted-count}
Let $X_1,\ldots,X_s\subseteq\mathbb P^N$ be varieties, $e\ge1$, and let
$h_1,\ldots,h_t$ be forms of degree at most $e$. Put
$Z=(\bigcup_kX_k)\cap V(h_1,\ldots,h_t)$, let $\mathcal I(Z)$ be its set of
isolated points and $\mathcal P(Z)$ its set of positive-dimensional
irreducible components. Then
\[
 \#\mathcal I(Z)+\sum_{Y\in\mathcal P(Z)}e^{\dim Y}\deg Y
 \le\sum_ke^{\dim X_k}\deg X_k .
\]
\end{lemma}

\begin{proof}
Assign to a variety $X$ the weight $e^{\dim X}\deg X$ and start with the list
$X_1,\ldots,X_s$. Process $h_1,\ldots,h_t$ in turn: keep every listed variety
contained in $V(h_i)$, and replace every other listed variety $X$ by the
irreducible components of $X\cap V(h_i)$; a point not in $V(h_i)$ is simply
removed. By (AG1), when $\dim X\ge1$ the new components have dimension
$\dim X-1$ and total degree at most $e\deg X$, so the total weight does not
increase. Keep repeated entries. After the last step the union of the listed
varieties is $Z$. Each isolated point of $Z$ and each $Y\in\mathcal P(Z)$ is a
listed variety: it is contained in some listed irreducible closed subset of
$Z$, which must coincide with it by maximality. These entries are distinct, so
their weights are bounded by the final, hence by the initial, total weight.
\end{proof}

\subsection{Degree bounds}
\label{app:algebraic-bounds}

\begin{lemma}[Flat directions]
\label{lem:algebraic-flat}
Let $f\in\Q[x_1,\ldots,x_n]$ be globally convex of degree at most four, with
$f\ge0$, a zero $p$, and $\nabla^2f(p)\succ0$, and let $f_4$, $f_3$ be its
homogeneous parts of degrees four and three. Then
$K_4=\{d:f_4(d)=0\}$ is a linear subspace defined over $\Q$, and $f_4$ and
$f_3$ are invariant under translations by $K_4$. In rational coordinates
$x=U(y,\vartheta)$, $U\in GL_n(\Q)$, with $K_4=\{y=0\}$, there are a rational
matrix $A_\vartheta\succ0$, a rational affine map $\vartheta_*$, and a
rational globally convex $\hat f$ of degree at most four whose quartic form is
positive definite, such that
\[
 f(U(y,\vartheta))=\hat f(y)+(\vartheta-\vartheta_*(y))^{\mathsf T}A_\vartheta
 (\vartheta-\vartheta_*(y)).
\]
The unique zero $\hat p$ of $\hat f$ is the $y$-part of $U^{-1}p$,
$\nabla^2\hat f(\hat p)\succ0$, and $\Q(\hat p)=\Q(p)$. If $f=\sum_jr_j^2$ with
rational quadratics $r_j$, then $\hat f=\sum_jr_j(U(y,\vartheta_*(y)))^2$ is a
sum of squares of rational quadratics.
\end{lemma}

\begin{proof}
As in Step 2 of the proof of Theorem~\ref{thm:sos-length}, $f_4$ is convex,
nonnegative and even, $K_4$ is a linear subspace, and $f_4(x+d)=f_4(x)$ for
$d\in K_4$. Hence $K_4=\{d:D_df_4\equiv0\}$, the kernel of a rational linear
map, and it has a rational basis. For $d\in K_4$,
$\nabla^2f(x+sd)=\nabla^2f(x)+s\nabla^2(D_df_3)$, since $\nabla^2f_4$ is
invariant and $\nabla^2f_3$ is linear in the point. Positive semidefiniteness
for all real $s$ forces $\nabla^2(D_df_3)=0$; the quadratic form $D_df_3$
therefore vanishes, and $f_3$ is invariant along $K_4$. In the new coordinates
the parts $f_4$, $f_3$ depend only on $y$, so
$f(U(y,\vartheta))=G(y)+\vartheta^{\mathsf T}A_\vartheta\vartheta
+\vartheta^{\mathsf T}(B_\vartheta y+c_\vartheta)$ with rational data, and
$A_\vartheta\succ0$ because it is half the $\vartheta\vartheta$ block of the
Hessian at $p$. Completing the square with
$\vartheta_*(y)=-\tfrac12A_\vartheta^{-1}(B_\vartheta y+c_\vartheta)$ gives the
displayed identity with $\hat f(y)=f(U(y,\vartheta_*(y)))$, which is convex as
the restriction of $f$ to an affine graph. Its quartic form is $f_4(U(y,0))$,
positive for $y\ne0$. Zeros of $\hat f$ correspond to zeros of $f$, and the
coordinates of $p$ are rational affine functions of those of $\hat p$ and
conversely. The Hessian of $\hat f$ at $\hat p$ is a Schur complement of the
positive definite Hessian of $f\circ U$. The last statement follows by
substitution.
\end{proof}

\begin{proof}[Proof of Theorem~\ref{thm:algebraic-degree-bounds}]
By Lemma~\ref{lem:one-real-conjugate} (zero-set case), $p$ is algebraic and
$D$ is odd. Write $R_j(x_0,x)=x_0^2r_j(x/x_0)$ for the quadratic forms
homogenizing the $r_j$, and identify $x\in\C^n$ with $[1:x]\in\mathbb P^n$.
Let $\Gamma$ be the set of the $D$ conjugates $\sigma(p)$ of $p$.

\emph{(a)} All $r_j$ vanish at $p$, and $\nabla^2f(p)=2Dr(p)^{\mathsf
T}Dr(p)$, so some $n$ of them, say $h_1,\ldots,h_n$, have nonsingular Jacobian
at $p$. Applying $\sigma$ to the rational determinant shows that their
Jacobian is nonsingular at every point of $\Gamma$. So every point of $\Gamma$
is a simple, hence isolated, zero of $h_1,\ldots,h_n$, and an isolated point of
$Z=V(H_1,\ldots,H_n)\subseteq\mathbb P^n$, where $H_i$ homogenizes $h_i$.
Lemma~\ref{lem:algebraic-weighted-count} with $X_1=\mathbb P^n$ and $e=2$ gives
$D\le2^n$, and oddness gives $D\le2^n-1$.

\emph{Reduction for (b) and (c).} Assume $f$ globally convex. Apply
Lemma~\ref{lem:algebraic-flat}. If $K_4\ne0$, then $\hat f$ is a rational sum of
quadratic squares in at most $n-1$ variables with the same field degree, so
part (a) gives $D\le2^{n-1}-1$, which is smaller than every bound in (b) and
(c); if no variable remains, $p$ is rational. So we may assume that $f_4$ is
positive definite. Then the forms $R_j$ have no common real zero at infinity:
$R_j(0,d)=r_{j,2}(d)$ for the quadratic parts, and $f_4(d)=\sum_jr_{j,2}(d)^2>0$
for real $d\ne0$. Their only common real affine zero is $p$.

\emph{(b), first bound.} Let $n\ge3$ and suppose $D=2^n-1$. With $Z$ as in (a),
Lemma~\ref{lem:algebraic-weighted-count} gives
$D+2\#\mathcal P(Z)\le2^n$, so $Z$ is finite. By (AG2) it is a complete
intersection of length $2^n$. Points of $\Gamma$ have local length one, so the
residual has length one: it is a single reduced point $b$, and $Z$ is
reduced. Since $Z$ and $\Gamma$ are defined over $\Q$, $b$ is fixed by
$\operatorname{Gal}(\overline\Q/\Q)$, so $b\in\mathbb P^n(\Q)$. Choose a
rational linear form $\Lambda_b$ with $\Lambda_b(b)\ne0$. For every $j$, the
form $R_j\Lambda_b^{\,n-3}$ has degree $n-1=\sum_i2-n-1$ and vanishes on
$\Gamma=Z\setminus\{b\}$, so by (AG3) it vanishes at $b$, hence $R_j(b)=0$. If
$b=[1:b']$, then $b'\in\Q^n$ is a real zero of $f$ different from $p$ (as
$b\notin\Gamma$). If $b=[0:d]$, then $f_4(d)=0$ with $d\ne0$ real. Both are
impossible. So $D\le2^n-3$.

\emph{(b), second bound.} Let $n\ge4$ and suppose $D=2^n-3$. Let
$\mathcal W=V(R_1,\ldots,R_m)\subseteq\mathbb P^n$. Its only real point is
$[1:p]$, which is isolated because the Jacobian of the $r_j$ has rank $n$.
Hence no positive-dimensional component of $\mathcal W$ has a real point.

\emph{Generic combinations.} For a matrix $\mathsf A\in\C^{n\times m}$ put
$h^{\mathsf A}=\mathsf Ar$ and $Z^{\mathsf A}=V(\mathsf AR)$. On
$U_{\mathcal W}=\mathbb P^n\setminus\mathcal W$ the vector $R(x)$ is nonzero,
so $\{(x,\mathsf A)\in U_{\mathcal W}\times\C^{n\times m}:\mathsf AR(x)=0\}$
is a vector bundle of rank $nm-n$ over the $n$-dimensional $U_{\mathcal W}$,
of dimension $nm$. By the theorem on the dimension of fibres, for
$\mathsf A$ in a nonempty Zariski open set the fibre $Z^{\mathsf A}\cap
U_{\mathcal W}$ is finite, so every positive-dimensional component of
$Z^{\mathsf A}$ lies in $\mathcal W$. The condition that $\mathsf A\,Dr(\sigma(p))$
be nonsingular for every $\sigma$ is another nonempty Zariski open condition.
A nonempty Zariski open subset of $\C^{n\times m}$ contains rational points,
so we may choose $\mathsf A$ rational with both properties. Put
$Z=Z^{\mathsf A}$; it is defined over $\Q$, and $\Gamma$ consists of simple
isolated points of $Z$.

\emph{Finiteness.} The set $\mathcal P(Z)$ is permuted by complex
conjugation, and its members lie in $\mathcal W$, so they have no real points.
If $\mathcal P(Z)\ne\emptyset$, its total weight
$\sum_{Y\in\mathcal P(Z)}2^{\dim Y}\deg Y$ is at least four: weight below four
would mean a single member, a line, which would then be conjugation-invariant
and have real points. Lemma~\ref{lem:algebraic-weighted-count} would give
$D+4\le2^n$, a contradiction. So $Z$ is finite, a complete intersection of
length $2^n$ by (AG2), and $Z=\Gamma\sqcup R$ with a residual subscheme $R$ of
length three, defined over $\Q$; $R$ is disjoint from $\Gamma$ because
points of $\Gamma$ have local length one.

\emph{The residual imposes independent conditions.} If $H_R(1)\le2$, the
linear forms vanishing on $R$ cut out a linear space of dimension at most one
containing $R$. A reduced point cannot contain a scheme of length three, so
$R$ lies on a line $\ell$. Each $\mathsf AR$ restricts to a binary quadratic
form on $\ell$ vanishing on a scheme of length three, hence vanishes on
$\ell$; then $\ell\subseteq Z$, contradicting finiteness. So $H_R(1)=3$. A
linear form $\Lambda$ that does not vanish at the support of $R$ is a unit on
$R$; multiplication by it maps the image of $S_k$ into that of $S_{k+1}$
injectively, so $H_R(k)=3$ for all $k\ge1$.

\emph{Conclusion.} Apply (AG3) with $s=n-1$ and $t=2$: since $n-3\ge1$,
$h^0(\mathcal I_\Gamma(2))-h^0(\mathcal I_Z(2))=3-H_R(n-3)=0$. Every $R_j$
vanishes on $\Gamma$, hence on $Z\supseteq R$. Since $R$ is defined over $\R$
and has odd length, it has a real point, which is a common zero of all $R_j$
outside $\Gamma$. As shown in the reduction, this is impossible. With oddness,
$D\le2^n-5$.

\emph{(c)} is proved in Appendix~\ref{app:algebraic-five}.

\emph{Maxima.} For $n=1$, (a) gives $D=1$. For $n=2,3,4,5$ the bounds are
$3,5,11,21$, and the quartics of Theorem~\ref{thm:algebraic-cyclic} are
globally convex rational sums of quadratic squares with field degrees
$d_2,\ldots,d_5=3,5,11,21$.
\end{proof}

\begin{proof}[Proof of Proposition~\ref{prop:algebraic-node-bound}]
The real zero set of $f$ is $\{p\}$, so Lemma~\ref{lem:one-real-conjugate}
shows that $p$ is algebraic and $D=[\Q(p):\Q]$ is odd. Since $f\ge0$ attains
its minimum at $p$, $\nabla f(p)=0$, and $\nabla^2f(p)\succ0$. For every
conjugate $b=\sigma(p)$: $f(b)=0$, $\nabla f(b)=0$, and the complex symmetric
matrix $\mathcal H_b=\nabla^2f(b)=\sigma(\nabla^2f(p))$ is nonsingular.

\emph{Generic gradient combinations.} For each $b$, the set of vectors
$\varpi\in\C^n$ with $\varpi^{\mathsf T}\mathcal H_b^{-1}\varpi=0$ is a proper
quadric cone. Choose $\varpi\in\Q^n$ outside these finitely many cones, and a
rational $(n-1)\times n$ matrix $\mathsf A$ of rank $n-1$ with
$\ker\mathsf A=\R\varpi$. At each $b$, the polynomial map $\mathsf A\nabla f$
has derivative $\mathsf A\mathcal H_b$ of rank $n-1$, so near $b$ its zero set
is a smooth complex curve $\gamma_b(t)$ with $\gamma_b(0)=b$ and tangent
$v_b=\mathcal H_b^{-1}\varpi$ (because $\mathsf A\mathcal H_bv_b=\mathsf
A\varpi=0$). The function $\varphi_b(t)=f(\gamma_b(t))$ satisfies
$\varphi_b(0)=\varphi_b'(0)=0$ and
$\varphi_b''(0)=v_b^{\mathsf T}\mathcal H_bv_b=\varpi^{\mathsf T}\mathcal
H_b^{-1}\varpi\ne0$.

\emph{Splitting the double points.} For a sufficiently small $c\ne0$, Rouch\'e's
theorem gives two distinct zeros of $\varphi_b(t)-c$ near $t=0$, at which
$\varphi_b'\ne0$. At such a point $x=\gamma_b(t)$ the Jacobian of the system
$(f-c,\ \mathsf A\nabla f)$ is nonsingular: the rows of $\mathsf A\nabla^2f(x)$
span the conormal space of the curve, and $\nabla f(x)^{\mathsf
T}\gamma_b'(t)=\varphi_b'(t)\ne0$. Taking $c$ small for all $D$ conjugates at
once, the system has at least $2D$ distinct simple, hence isolated,
solutions.

\emph{B\'ezout.} Homogenize $f-c$ to a form of degree $\deg f\le4$ and the
entries of $\mathsf A\nabla f$ to forms of degree at most three. The
irreducible components of the quartic hypersurface have dimension $n-1$ and
total degree at most four, so their total weight with $e=3$ is at most
$4\cdot3^{n-1}$. Lemma~\ref{lem:algebraic-weighted-count} bounds the number of
isolated points of the full intersection, which include the $2D$ affine
isolated solutions, by $4\cdot3^{n-1}$. Hence $D\le2\cdot3^{n-1}$, and
oddness gives $D\le2\cdot3^{n-1}-1$.

\emph{Two variables with convexity.} Let $n=2$ and $f$ globally convex. A
convex polynomial of odd degree at least three does not exist, because its
leading form would be convex and odd, hence linear. If $\deg f\le2$, $p$ is
the minimizer of a rational strictly convex quadratic, so it is rational. If $f_4$ has a nonzero real zero, Lemma~\ref{lem:algebraic-flat}
reduces to one variable, where the bound just proved gives $D\le1$. Otherwise
$\deg f=4$ and $f_4$ is positive definite; Lemma~\ref{lem:algebraic-bivariate-descent}
shows that $f$ is a sum of squares of rational polynomials, necessarily
quadratics, and Theorem~\ref{thm:algebraic-degree-bounds}(a) gives $D\le3$.
\end{proof}

Lemma~\ref{lem:algebraic-bivariate-descent} uses the following classification.

\begin{enumerate}[label=(Sch)]
\item \cite{Scheiderer2016} Let $F\in\Q[x_0,x_1,x_2]$ be a quartic form that is
  nonnegative on $\R^3$ but is not a sum of squares of quadratic forms with
  rational coefficients. Then $F=\lambda\ell_1\ell_2\ell_3\ell_4$ with
  $\lambda\in\C$ and linear forms $\ell_i\in\C[x_0,x_1,x_2]$ such that no three
  of the lines $V(\ell_i)$ pass through a common point.
\end{enumerate}

\begin{proof}[Proof of Lemma~\ref{lem:algebraic-bivariate-descent}]
Let $F(x_0,x)=x_0^4f(x/x_0)$. Then $F\ge0$ on $\R^3$: for $x_0\ne0$ by
homogeneity, and for $x_0=0$ because $F(0,x)=f_4(x)$. Its real projective
zeros are the real zero of $f$ only, since $f_4(x)>0$ for real $x\ne0$.
Suppose $F$ is not a sum of squares of rational quadratic forms, and write
$F=\lambda\ell_1\ell_2\ell_3\ell_4$ as in (Sch). The four lines are distinct.
No $\ell_i$ is proportional to a real form: otherwise $F=\ell_iK$ with a real
cubic form $K$ not vanishing identically on the real line $V(\ell_i)$, and
near a real point of that line where $K\ne0$, $F$ would change sign. Since
$F$ is real and its factorization into lines is unique, conjugation permutes
the four lines without fixed points; they form two conjugate pairs
$\{\ell_1,\overline{\ell_1}\}$ and $\{\ell_2,\overline{\ell_2}\}$. The point
$V(\ell_1)\cap V(\overline{\ell_1})$ is conjugation invariant, hence real, and
it is a zero of $F$; so is $V(\ell_2)\cap V(\overline{\ell_2})$, and the two
points differ because no three lines are concurrent. This gives two real
projective zeros, a contradiction. Hence $F=\sum_iG_i^2$ with rational
quadratic forms $G_i$, and $f(x)=\sum_iG_i(1,x)^2$.

For a certified quartic $(f,A)$ in two variables with minimum zero, the full
Hessian Gram gives $12f_4(x)=(x\otimes x)^{\mathsf T}A_{22}(x\otimes x)>0$ for
$x\ne0$, where $A_{22}$ is the lower right block of $A$, and strong convexity
gives a unique zero; so the hypotheses hold.
\end{proof}

\begin{example}[Global convexity is needed in Theorem~\ref{thm:algebraic-degree-bounds}(b)]
\label{ex:algebraic-seven}
Let $r_1=z-x^2$, $r_2=yz-y^2-1$, $r_3=xz+y^2$ and $f=r_1^2+r_2^2+r_3^2$ on
$\R^3$. A common zero has $z=x^2$ and $y^2=-x^3$, and then $r_2=0$ reads
$yx^2+x^3-1=0$; $x=0$ is impossible, so $y=(1-x^3)/x^2$ and
$P_7(x)=x^7+x^6-2x^3+1=0$. Conversely every root of $P_7$ gives a common zero.
Modulo two, $\bar P_7=x^7+x^6+1$ has no root in $\mathbb F_2$, and repeated
squaring in $\mathbb F_2[x]/(\bar P_7)$, using
$x^8\equiv x^6+x+1$, $x^{10}\equiv x^6+x^3+x^2+x+1$, and
$x^{12}\equiv x^6+x^5+x^4+x^3+x^2+x+1$, gives
\[
 x^{16}\equiv x^6+x^5+x^4+x^3+x,\quad
 x^{32}\equiv x^5+x^4+x^2+x+1,\quad
 x^{64}\equiv x^4+x^3+1,\quad
 x^{128}\equiv x .
\]
Since $7$ is prime, $\bar P_7$ is irreducible, hence so is $P_7$ over $\Q$.
For $x\ge0$, $P_7(x)=x^7+(x^3-1)^2>0$; $P_7(-t)=-t^7+t^6+2t^3+1$ has one sign
change, so by Descartes' rule $P_7$ has at most one negative root, and
$P_7(-2)<0<P_7(-1)$. Thus $f$ has exactly one real zero $p$, with
$\Q(p)=\Q(x)$ of degree $7>2^3-3$. At every common zero the Jacobian
determinant of $(r_1,r_2,r_3)$ equals $-P_7'(x)/x^2\ne0$, so
$\nabla^2f(p)=2Dr^{\mathsf T}Dr\succ0$. However
$\partial_x^2f(0,0,1)=-2$, and the homogenized quadrics share the real point
$[x:y:z:x_0]=[0:0:1:0]$ at infinity.
\end{example}

\subsection{Five variables}
\label{app:algebraic-five}

\begin{proposition}[Finite odd residuals in $\mathbb P^5$]
\label{prop:algebraic-residual}
Let $Z\subset\mathbb P^5$ be a zero-dimensional complete intersection of five
quadrics with rational coefficients, and let $Z=\Gamma\sqcup R$ be a
decomposition into closed subschemes defined over $\Q$, with $\Gamma$ reduced
and $\operatorname{length}(R)\in\{1,3,5,7,9\}$. Then some real point of $R$ is
a common zero of all quadratic forms that vanish on $\Gamma$.
\end{proposition}

\begin{proof}
Choose a rational linear form $\Lambda$ that does not vanish at any point of
$Z$, and work in the affine chart $\Lambda=1$. Let $\mathcal A$ be the
coordinate ring of $Z$ over $\Q$ in this chart, a $\Q$-algebra of dimension
$32$, filtered by the images $F_i\mathcal A$ of polynomials of degree at most
$i$. By (AG2), $\Lambda$ is a nonzerodivisor on the homogeneous coordinate ring
$S_\Q/I_Z$, so $F_i\mathcal A\cong(S_\Q/I_Z)_i$ and the associated graded algebra
is $S_\Q/(I_Z+(\Lambda))$, a complete intersection of five quadrics in five
variables. By (AG4) it has Hilbert series $(1+t)^5$ and perfect pairings
$\operatorname{gr}_i\times\operatorname{gr}_{5-i}\to\operatorname{gr}_5$.

\emph{A nondegenerate functional.} Let $\lambda:\mathcal A\to\Q$ vanish on
$F_4\mathcal A$ and be nonzero on the one-dimensional space
$\mathcal A/F_4\mathcal A=\operatorname{gr}_5$. If $a\ne0$ has leading degree
$i$, choose $b'\in\operatorname{gr}_{5-i}$ pairing nontrivially with its leading
class and lift it to $b\in F_{5-i}\mathcal A$; then $ab\in F_5\mathcal A$ has
nonzero image in $\operatorname{gr}_5$, so $\lambda(ab)\ne0$. Hence
$(a,b)\mapsto\lambda(ab)$ is nondegenerate on $\mathcal A$. The decomposition
$Z=\Gamma\sqcup R$ gives $\mathcal A=\mathcal A_\Gamma\times\mathcal A_R$, and
the restriction $\lambda_R$ of $\lambda$ to $\mathcal A_R$ is again
nondegenerate.

\emph{The quotient by an annihilator.} Let $\mathsf q$ be a rational quadric
vanishing on $\Gamma$, dehomogenized in the chart; its image in $\mathcal A$ is
$(0,\mathsf q_R)$. Put $\mathcal B=\mathcal A_R/\operatorname{Ann}(\mathsf q_R)$,
$\ell=\dim\mathcal B$, and $\beta(a,b)=\lambda_R(\mathsf q_Rab)$. This is a
well-defined nondegenerate symmetric bilinear form on $\mathcal B$: an element
pairs to zero with everything exactly when its product with $\mathsf q_R$
vanishes. Let $\mathcal V\subseteq\mathcal B$ be the image of the polynomials
of degree at most one and $v=\dim\mathcal V$. For $a,b$ of degree at most one,
$\mathsf qab$ has degree at most four and vanishes on $\Gamma$, so
$\beta(a,b)=\lambda(\mathsf qab)=0$. Thus $\mathcal V$ is totally isotropic
and $2v\le\ell$.

\emph{A B\'ezout bound.} The subscheme $T=\operatorname{Spec}\mathcal B$ of $R$
lies in the projective linear space $\Pi\cong\mathbb P^{v-1}$ cut out by the
linear forms vanishing on it. The five quadrics restrict to quadrics on $\Pi$
whose common zero set $Z\cap\Pi$ is finite; choosing successively general
combinations that do not contain any positive-dimensional component of the
current intersection gives $v-1$ quadrics on $\Pi$ with finite common zero
scheme containing $T$. By (AG2) its length is $2^{v-1}$, so
$\ell\le2^{v-1}$ (for $v=1$, $T$ is contained in a reduced point and
$\ell\le1$). If $1\le\ell\le9$, the inequalities $2v\le\ell\le2^{v-1}$ force
$v=4$ and $\ell=8$.

\emph{Real parity.} Suppose that no real point of $R$ is a common zero of all
quadrics through $\Gamma$. Then a suitable rational combination $\mathsf q$ of
a basis of these quadrics is nonzero at every one of the finitely many real
points of $R$. Over $\R$, $\mathcal A_R\otimes\R$ is a product of local
algebras with residue fields $\R$ (at real points) or $\C$ (at pairs of
conjugate points). On a factor with residue field $\R$, $\mathsf q$ is a unit,
so the factor survives unchanged in $\mathcal B\otimes\R$. A factor with
residue field $\C$ is a $\C$-algebra, so it and its quotients have even real
dimension. Hence $\ell\equiv\operatorname{length}(R)\pmod2$ is odd. Since
$\operatorname{length}(R)$ is odd, $R$ has a real point, $\mathsf q$ is nonzero
there, and $1\le\ell\le9$. This contradicts $\ell=8$.
\end{proof}

\begin{proposition}[Positive-dimensional base loci in $\mathbb P^5$]
\label{prop:algebraic-positive-base}
Let $\mathcal Q$ be a complex vector space of quadratic forms on
$\mathbb P^5$ spanned by real forms, with base locus $\mathcal W$. Suppose no
positive-dimensional component of $\mathcal W$ has a real point, and that
$\Gamma\subset\mathcal W$ is a set of $D\ge23$ points at each of which the
differentials of the forms in $\mathcal Q$ span a space of dimension five
(in an affine chart). Then $\mathcal W$ is finite.
\end{proposition}

\begin{proof}
Suppose $\mathcal W$ has positive-dimensional components. Each point of
$\Gamma$ is an isolated point of $\mathcal W$, so it lies on none of them.

\emph{Weight.} As in the proof of Theorem~\ref{thm:algebraic-degree-bounds},
five general forms $h_1,\ldots,h_5\in\mathcal Q$ have simple common zeros at
$\Gamma$ and no positive-dimensional common component outside $\mathcal W$;
every positive-dimensional component of $\mathcal W$ is then a component of
$V(h_1,\ldots,h_5)$. Lemma~\ref{lem:algebraic-weighted-count} gives
\begin{equation}
 \sum_{Y}2^{\dim Y}\deg Y\le32-D\le9,
 \label{eq:algebraic-base-weight}
\end{equation}
the sum running over the positive-dimensional components of $\mathcal W$.
This set is permuted by conjugation.

\emph{The possible configurations.} Components of dimension at least four
have weight at least $16$. A three-dimensional component has weight at least
$8$, so it would be the only positive-dimensional component, conjugation
invariant, of degree one, hence a real linear space, which has real points.
If there is a surface component, the surface components have total degree at
most two: either a single plane, which would be invariant and real; two
planes, which must be conjugate (invariant planes are real); or one integral
quadric surface, which is invariant. In the last two cases no curve component
fits in \eqref{eq:algebraic-base-weight}. Otherwise all components are curves of total
degree at most four. An invariant curve has even degree (odd degree would give
a real point), so it is an integral conic or an integral quartic. A
non-invariant curve comes with its conjugate, and both have degree one or two.
In each case we choose a subvariety $Y_0\subseteq\mathcal W$ in the list below.
It is a nonempty reduced pure-dimensional local complete intersection with
$\mathcal I_{Y_0}(2)$ globally generated, and $e(Y_0)\ge10$.

\emph{Two conjugate planes.} They are disjoint, since their intersection
would be an invariant linear space with real points. In coordinates with the
planes $\{x_3=x_4=x_5=0\}$ and $\{x_0=x_1=x_2=0\}$, the nine products
$x_ix_j$, $i\le2<j$, generate the ideal. For a plane, $N=\mathcal O(1)^3$ and
$e=[H^2](1+2H)^5(1+H)^{-3}=40-30+6=16$; so $e(Y_0)=32$.

\emph{A real integral quadric surface.} By (AG6) an irreducible surface of
degree two spans a linear space of dimension at most three; it cannot span a
plane, since then it would be the plane, of degree one. So it spans a
$\mathbb P^3$, which is real because the surface is conjugation invariant, and
it is a hypersurface of degree two there. An irreducible and reduced quadric in
$\mathbb P^3$ has rank three or four; if it had rank three, its unique
singular point would be conjugation invariant, hence a real point of
$\mathcal W$ on a positive-dimensional component. So it has rank four and is
smooth. Its ideal is generated by two linear forms and
one quadric, $N=\mathcal O(1)^2\oplus\mathcal O(2)$, and
$e=2[H^2](1+2H)^4(1+H)^{-2}=2(24-16+3)=22$.

\emph{A real integral conic.} By (AG6) it spans a plane, in which it is a
smooth conic, cut out by three linear forms and a quadric; (AG7) gives
$e=4\cdot2+2=10$.

\emph{A real integral quartic curve $C$.} By (AG6) its span has dimension at
most four. If the span is a plane, every quadric of $\mathcal Q$ restricts to
a conic containing the integral plane quartic, hence vanishes on the plane;
the plane would lie in $\mathcal W$, contradicting the configuration. If the
span is a $\mathbb P^3$ (necessarily real) and the quadrics of $\mathcal Q$
restrict there to a space of dimension at most one, $\mathcal W$ contains the
span or a quadric surface, again a contradiction. Otherwise two independent
restricted quadrics contain $C$; they have no common component, since a common
plane would contain the nondegenerate curve or leave it on a line. Their
complete intersection has degree four and contains the reduced curve $C$ of
degree four; being unmixed and generically reduced, it equals $C$. So $C$ is a
complete intersection of type $(1,1,2,2)$ in $\mathbb P^5$, with
$N=\mathcal O(1)^2\oplus\mathcal O(2)^2$ and
$e=4[H](1+2H)^3(1+H)^{-2}=4(6-2)=16$, smooth or not. If the span is
$\mathbb P^4$, $C$ is a rational normal quartic by (AG6), with ideal generated
by quadrics and one linear form, and (AG7) gives $e=4\cdot4+2=18$.

\emph{Two conjugate lines.} They are disjoint, span a real $\mathbb P^3$, and
in coordinates there with lines $\{x_0=x_1=0\}$, $\{x_2=x_3=0\}$ the four
products $x_ix_j$, $i\le1<j$, together with two linear forms generate the
ideal. Each line has $e=6$, so $e(Y_0)=12$.

\emph{Two conjugate conics $C,\overline C$ in planes $\Pi,\overline\Pi$.} If
$\Pi=\overline\Pi$, the plane is real, and every quadric of $\mathcal Q$
restricts to a conic vanishing on two distinct integral conics, hence vanishes
on the plane; this contradicts the configuration. If the planes are disjoint,
the products between the two coordinate blocks together with the two conic
equations generate the ideal sheaf of $C\sqcup\overline C$; then
$e=10+10=20$. If the planes meet in one point $a$, it is real and lies on
neither conic (else it would be a real point of $\mathcal W$ on a curve
component). In coordinates on the span $\mathbb P^4$ with $\Pi=V(x_3,x_4)$,
$\overline\Pi=V(x_1,x_2)$, $a=[1:0:0:0:0]$, and conic equations
$c(x_0,x_1,x_2)$, $c'(x_0,x_3,x_4)$ normalized to have coefficient one on
$x_0^2$, the four products $x_ix_j$ ($i\in\{1,2\}$, $j\in\{3,4\}$), the
quadric $c+c'-x_0^2$, and the linear form of the span generate the ideal sheaf
of $C\sqcup\overline C$: away from $a$ the products cut out one plane locally
and the quadric restricts to that plane's conic, while at $a$ the quadric
does not vanish. Again $e=20$. If the planes meet along a line, their span is
a real $\mathbb P^3$ and $\Pi\cup\overline\Pi=V(\mathsf Q_0)$ for a real
quadric $\mathsf Q_0$ there. Some real $\mathsf q\in\mathcal Q$ does not
vanish on $\Pi$ (otherwise, by conjugation, $\mathcal W$ would contain both
planes and their real common line). Its restrictions to $\Pi$ and
$\overline\Pi$ are nonzero multiples of the two conic equations, so
$V(\mathsf Q_0,\mathsf q)$ in that $\mathbb P^3$ is a complete intersection of
degree four containing $C\cup\overline C$; as before it equals
$C\cup\overline C$, of type $(1,1,2,2)$, and $e=16$.

\emph{Conclusion.} In every case $Y_0\subseteq\mathcal W$, so the five forms
$h_i$ vanish on $Y_0$ and have the $D$ simple zeros $\Gamma$ outside $Y_0$. By
the implicit function theorem, every quintuple of quadrics vanishing on $Y_0$
that is close enough to $(h_1,\ldots,h_5)$ has at least $D$ distinct simple
zeros outside $Y_0$, one near each point of $\Gamma$. General quintuples in
$H^0(\mathcal I_{Y_0}(2))^5$ are dense, so by (AG5) $D\le32-e(Y_0)\le22$, a
contradiction.
\end{proof}

\begin{proof}[Proof of Theorem~\ref{thm:algebraic-degree-bounds}(c)]
Let $n=5$, $f$ globally convex, and suppose $D\ge23$; by (a) and oddness,
$D\in\{23,25,27,29,31\}$. As in the reduction of the proof of (b), we may
assume that $f_4$ is positive definite, since otherwise $D\le15$. Let
$\mathcal Q$ be the complex span of the forms $R_j$, $\mathcal W$ its base
locus, and $\Gamma$ the orbit of $p$. The only real point of $\mathcal W$ is
$[1:p]$, and it is isolated, so positive-dimensional components of
$\mathcal W$ have no real points; the differentials of $\mathcal Q$ span a
five-dimensional space at each point of $\Gamma$.
Proposition~\ref{prop:algebraic-positive-base} shows that $\mathcal W$ is
finite. Choose a rational $5\times m$ matrix $\mathsf A$ as in the proof of
(b); then $Z=V(\mathsf AR)$ has no positive-dimensional component (those would
lie in the finite $\mathcal W$), so it is a zero-dimensional complete
intersection of length $32$, defined over $\Q$, with the simple points
$\Gamma$. Its residual $R$ has length $32-D\in\{1,3,5,7,9\}$ and is defined
over $\Q$. Proposition~\ref{prop:algebraic-residual} gives a real point of $R$
at which every $R_j$ vanishes; it is not in $\Gamma$. As in (b), this is a
second real zero of $f$ or a real zero of $f_4$ at infinity, a contradiction.
Hence $D\le21$.
\end{proof}

\subsection{Univariate realizations}
\label{app:algebraic-univariate}

For a univariate $f$ and a centre $c$, put $\mathsf z=x-c$ and
$\mathsf w=(1+\mathsf z^2)^N$. Direct differentiation of
$g=f\mathsf w$ gives the identity
\begin{equation}
 \frac{g''}{\mathsf w}=f''+\frac{4N\mathsf z}{1+\mathsf z^2}f'
 +\Bigl[\frac{2N}{1+\mathsf z^2}+\frac{4N(N-1)\mathsf z^2}{(1+\mathsf z^2)^2}
 \Bigr]f .
 \label{eq:algebraic-multiplier}
\end{equation}

\begin{proof}[Proof of Theorem~\ref{thm:algebraic-univariate}(b)]
Write $f=P^2$, of degree $2d$, and define the rational constants
\[
\begin{gathered}
 R_0=2^{\tau+1},\quad \sigma_1=2^{-4d^2(\tau+1)},\quad
 m_1=\sigma_1^{2(d-1)},\quad C_0=(d+1)2^{2\tau},\quad
 Q_0=(2d+1)^2C_0,\\
 R_1=4(R_0+Q_0+1),\quad H_1=(2d+1)^4C_0R_1^{2d},\quad
 r_1=\frac{m_1}{2H_1},\quad \phi_1=r_1^2(\sigma_1/2)^{2(d-1)},\\
 \psi_1=\frac{\phi_1r_1^2}{(1+4R_1^2)^2},\quad
 N=2+\Bigl\lceil\frac{3H_1+1}{\psi_1}\Bigr\rceil,\quad \kappa=\frac2{m_1}.
\end{gathered}
\]
\emph{Bounds on $f$.} All roots satisfy $\abs\gamma<R_0$. As in the proof of
Theorem~\ref{thm:algebraic-dimension}, the nonzero integer discriminant gives
$\abs{\gamma-\gamma'}\ge2^{-\tau(d-1)}(2R_0)^{-d(d-1)/2+1}>\sigma_1$ for distinct
roots when $d\ge2$. Hence $\abs{P'(\alpha)}\ge\sigma_1^{d-1}$ (the leading
coefficient is a positive integer) and $f''(\alpha)=2P'(\alpha)^2\ge2m_1$. The
coefficients of $f$ are bounded by $C_0$ and those of $f',f''$ by $Q_0$, so
$\abs{f^{(j)}(x)}\le H_1$ for $\abs x\le R_1$ and $j\le3$. We use three regions.
\begin{enumerate}[label=(\Roman*)]
\item If $\abs{x-\alpha}\le r_1$, then $\abs x<R_1$, $f''(x)\ge2m_1-H_1r_1\ge
  m_1$, $f''\le H_1$, and $f'$ has the sign of $x-\alpha$ with
  $\abs{f'(x)}\le H_1\abs{x-\alpha}$.
\item If $\abs x\le R_1$ and $\abs{x-\alpha}\ge r_1$, then
  $f(x)=a_d^2\prod_\gamma\abs{x-\gamma}^2\ge r_1^2(\sigma_1/2)^{2(d-1)}=\phi_1$,
  because $\abs{x-\gamma}\ge\abs{\operatorname{Im}\gamma}>\sigma_1/2$ for the
  nonreal roots.
\item If $\abs x\ge R_1$, then $f''(x)\ge1$ and $xf'(x)>0$: the leading
  coefficients of $f'$ and $f''$ are at least two, and since $R_1>4(Q_0+1)$
  the lower terms are less than a quarter of the leading term in absolute
  value. (For $d=1$, $f''$ is a constant at least two.)
\end{enumerate}
\emph{The centre.} Because $P$ has a single, simple real root and a positive
leading coefficient, $d$ is odd and $P(-R_0)<0<P(R_0)$. Rational bisection
with exact sign evaluation by Horner's rule produces a dyadic $c$ with
$\delta=\abs{c-\alpha}\le r_1/2$ and $\delta^2\le m_1/(8NH_1)$ after a number of
steps polynomial in $d+\tau$; if a midpoint is a root, take it.

\emph{Curvature.} In region (I), the last term of
\eqref{eq:algebraic-multiplier} is nonnegative, and the middle term is
negative only when $x$ lies between $c$ and $\alpha$, where
$\abs{\mathsf z}\le\delta$ and $\abs{f'}\le H_1\delta$; so
$g''/\mathsf w\ge m_1-4NH_1\delta^2\ge m_1/2$. In region (II),
$r_1/2\le\abs{\mathsf z}\le2R_1$, so the last term is at least
$N(N-1)\phi_1r_1^2/(1+4R_1^2)^2=N(N-1)\psi_1$, the middle term is at least
$-2NH_1$ (as $\abs{\mathsf z}/(1+\mathsf z^2)\le1/2$), and
$g''/\mathsf w\ge N(N-1)\psi_1-2NH_1-H_1\ge N(H_1+1)-H_1\ge1$ because
$(N-1)\psi_1\ge3H_1+1$. In region (III), $\abs c<R_0+1<R_1\le\abs x$, so
$\mathsf z$ and $x$ have the same sign, $\mathsf zf'>0$, and
$g''/\mathsf w\ge f''\ge1$. Since $\mathsf w\ge1$ and $m_1\le1$,
$g''\ge m_1/2$ everywhere, and $G=\kappa g$ satisfies $G''\ge1$. The factor
$\mathsf w$ is positive, so $G^{-1}(0)=\{\alpha\}$.

\emph{Size.} The logarithms of $R_0,R_1,H_1$ are $O(d(\tau+\log(d+1)))$, and
those of $1/\sigma_1$, $1/m_1$, $1/r_1$, $1/\phi_1$, $1/\psi_1$, $N$ and $\kappa$
are $O(d^3(\tau+\log(d+1)))$. A circuit computes $P$ by Horner's rule, squares
it, forms $1+(T-c)^2$, raises it to the power $N$ by repeated squaring with
shared intermediate results ($O(\log N)$ gates), and multiplies by $\kappa$.
For the integer $x_0=2^{\tau+2}$, $x_0-c\ge2^{\tau+1}\ge4$ and $P(x_0)$ is a
nonzero integer, so $G(x_0)\ge\kappa2^N\ge2^N$; its reduced numerator has at
least $N+1$ bits.
\end{proof}

\begin{proof}[Proof of Theorem~\ref{thm:algebraic-univariate}(a)]
If $\alpha$ is the unique real zero of a rational univariate polynomial,
Lemma~\ref{lem:one-real-conjugate} (zero-set case, $n=1$) shows that it is
algebraic with one real conjugate. Conversely, apply part (b) to the primitive
integer multiple of the minimal polynomial of $\alpha$.
\end{proof}

The order of choices in part (b) matters: the exponent $N$ is fixed first and
the centre is then chosen close to $\alpha$. A fixed centre does not suffice in
general. For $f=(x-1)^2$, $c=0$ and $x_N=1-2/N$, the right side of
\eqref{eq:algebraic-multiplier} tends to $2-8+4=-2$ as $N\to\infty$, so
$f(x)(1+x^2)^N$ is not convex for large $N$.

\begin{proof}[Proof of Theorem~\ref{thm:algebraic-univariate}(c)]
\emph{The cubic.} $P_k'(x)=3(x^2-1)$, $P_k(-1)=4+2^{-2k}>0$ and
$P_k(1)=2^{-2k}>0$, so $P_k$ has exactly one real root, and $P_k(-2)=2^{-2k}>0$
shows that it is less than $-2$. Write $\alpha_k=-2-t$ with $t>0$; then
$P_k(\alpha_k)=0$ reads $t(3+t)^2=2^{-2k}$. If $t=a/b$ in lowest terms with
positive integers, then $t(3+t)^2=a(a+3b)^2/b^3$ in lowest terms (both $a$ and
$a+3b$ are coprime to $b$), with numerator at least $16$; but the reduced
numerator of $2^{-2k}$ is $1$. So $P_k$ has no rational root and, being a
cubic, is irreducible. Comparing coefficients, the other roots are
$\beta\pm\mathrm i\eta$ with
\[
 \beta=1+\tfrac t2,\qquad \eta^2=3t+\tfrac34t^2,\qquad
 \beta-\alpha_k=3+\tfrac32t>3,\qquad 2^{-2k}-\eta^2=t\bigl(6+\tfrac{21}4t+t^2\bigr)>0,
\]
so $0<\eta<2^{-k}$. Clearing denominators shows that $P_k$ has binary length
$O(k)$.

\emph{An interval estimate.} We use Markov's inequality: a real polynomial
$\mathcal P$ of degree at most $D_1$ satisfies
$\max_{[a,b]}\abs{\mathcal P'}\le\frac{2D_1^2}{b-a}\max_{[a,b]}\abs{\mathcal P}$
\cite[Theorem~5.1.8]{BorweinErdelyi1995}. Iterating it with the same $D_1$
bounds the $j$-th derivative by $(2D_1^2/(b-a))^j\max\abs{\mathcal P}$. Suppose
$\mathcal P(b)=1$, $0\le\mathcal P\le1$ on $[a,b]$, and
$\mathcal P(b+\mathrm i\eta)=0$. Taylor expansion at $b$ gives
$1\le\sum_{j=1}^{D_1}q_M^j/j!$ with $q_M=2D_1^2\abs\eta/(b-a)$. If $q_M<1/2$ the
right side is less than $q_M/(1-q_M)<1$; hence
\begin{equation}
 D_1^2\ge\frac{b-a}{4\abs\eta}.
 \label{eq:algebraic-markov}
\end{equation}

\emph{Convex realizations.} Let $F$ be rational, globally convex, with only
real zero $\alpha_k$. Since $F(\alpha_k)=0$ and $P_k$ is irreducible,
$P_k\mid F$, so $F(\beta+\mathrm i\eta)=0$; in particular $F$ is not affine.
A nonaffine convex polynomial tends to $+\infty$ at both ends, so a negative
value would produce two real zeros; hence $F\ge0$ and $F(\beta)>0$. On
$[\alpha_k,\beta]$, $\mathcal P=F/F(\beta)$ lies between $0$ and the chord
from $0$ to $1$, so \eqref{eq:algebraic-markov} with $a=\alpha_k$, $b=\beta$
gives $(\deg F)^2\ge(\beta-\alpha_k)/(4\eta)>\tfrac342^k$. The normalizing
constant may be irrational; rationality is used only to force the complex
zero.

\emph{Systems.} Remove identically zero rows from a finite system of rational
convex inequalities $F_i\le0$ with solution set $\{\alpha_k\}$. Some active
row $F_i$ has $F_i'(\alpha_k)\ge0$: otherwise all active rows would be
negative on a short interval to the right of $\alpha_k$, where the inactive
rows are also negative. Its tangent line gives $F_i\ge0$ on
$[\alpha_k,\infty)$, and $F_i(\beta)>0$, since otherwise convexity would make
$F_i$ vanish on $[\alpha_k,\beta]$. As before $P_k\mid F_i$, and the same
argument applies to $F_i/F_i(\beta)$.

\emph{Objectives.} If a nonconstant rational convex $J$ has a minimizer at
$\alpha_k$, then $J'(\alpha_k)=0$, so $P_k\mid J'$. Convexity makes $J'$
nondecreasing, so $0\le J'\le J'(\beta)$ on $[\alpha_k,\beta]$ with
$J'(\beta)>0$ (otherwise $J'$ vanishes on an interval and $J$ is constant).
Apply \eqref{eq:algebraic-markov} to $J'/J'(\beta)$, of degree $\deg J-1$.
\end{proof}

\begin{proof}[Proof of Theorem~\ref{thm:algebraic-univariate}(d)]
Theorem~\ref{thm:algebraic-dimension} with $d=3$ and $n=2$ applies to $P_k$.
For $J_k$, $J_k'=P_k$ is negative left of $\alpha_k$ and positive right of it,
so $J_k$ decreases and then increases: its sublevel sets are intervals, and
$\alpha_k$ is its unique minimizer. But $J_k''(0)=-3$.
\end{proof}

\subsection{The quadratic graph}
\label{app:algebraic-graph}

\begin{proof}[Proof of Theorem~\ref{thm:quadratic-graph}]
Let $h=n+n^2$, $U_A=\tr A$, and $\varrho_A=\det A/U_A^{h-1}$. Every eigenvalue
of $A$ is at most $U_A$, so $A\succeq\varrho_AI$, and therefore
$\nabla^2f\succeq\varrho_AI$ on $\R^n$. Strong monotonicity of $\nabla f$ gives
$\varrho_A\norm p^2\le-\nabla f(0)^{\mathsf T}p$, so
$\norm p\le P_A:=1+\norm{\nabla f(0)}_1/\varrho_A$. All these rationals have
polynomial bit length.

\emph{Approximate minimizers.} For a rational $0<\eta\le1$, a rational
$\hat p$ with $\norm{\hat p-p}\le\eta$ is computed in time polynomial in $L$
and $\log(1/\eta)$: Lemma~\ref{lem:convex-value}, applied to $f$ on the box
$[-P_A-1,P_A+1]^n$, which contains $p$, returns a rational $\hat p$ in the box
with $f(\hat p)\le f(p)+\varrho_A\eta^2/2$, and strong convexity gives
$\frac{\varrho_A}2\norm{\hat p-p}^2\le f(\hat p)-f(p)$. The promise
$\min f=0$ is not used here.

\emph{A quadratic lift of $f$.} Let $C_{\mathrm{dup}}$ map
$(s_{ij})_{i\le j}$ to the symmetric $n\times n$ matrix, viewed as a vector
in $\R^{n^2}$; its singular values lie in $[1,\sqrt2]$. For rational $\hat p$
put $d=y-\hat p$ and $s_{ij}=z_{ij}-\hat p_iy_j-\hat p_jy_i+\hat p_i\hat p_j$,
and define
\[
 g_{\hat p}(y,z)=f(\hat p)+\nabla f(\hat p)^{\mathsf T}d
 +\int_0^1(1-\varsigma)
 \begin{pmatrix}d\\ \hat p\otimes d+\varsigma C_{\mathrm{dup}}s\end{pmatrix}^{\mathsf T}
 A\begin{pmatrix}d\\ \hat p\otimes d+\varsigma C_{\mathrm{dup}}s\end{pmatrix}
 d\varsigma .
\]
The integrand is a quadratic polynomial in $(y,z)$ whose coefficients are
polynomials of degree at most two in $\varsigma$, and
$\int_0^1(1-\varsigma)\varsigma^kd\varsigma=\tfrac12,\tfrac16,\tfrac1{12}$ for
$k=0,1,2$; so $g_{\hat p}$ is a rational quadratic. On the graph
$z_{ij}=y_iy_j$ we have $s_{ij}=d_id_j$, so $C_{\mathrm{dup}}s=d\otimes d$ and
$\hat p\otimes d+\varsigma d\otimes d=(\hat p+\varsigma d)\otimes d$. Taylor's
formula with integral remainder and the full Hessian Gram give
\[
 g_{\hat p}(y,(y_iy_j))=f(\hat p)+\nabla f(\hat p)^{\mathsf T}d
 +\int_0^1(1-\varsigma)\,d^{\mathsf T}\nabla^2f(\hat p+\varsigma d)\,d\,d\varsigma
 =f(y).
\]
In particular $g_{\hat p}(w_*)=f(p)=0$ for every rational $\hat p$.

\emph{Uniform bounds.} Assume $\norm{\hat p-p}\le1$, so
$\norm{\hat p}\le P_A+1$, and put $J_A=4(P_A+2)$. The quadratic part of
$g_{\hat p}$ is
$\int_0^1(1-\varsigma)(y,\hat p\otimes y+\varsigma C_{\mathrm{dup}}s^{\rm
lin})^{\mathsf T}A(\cdots)\,d\varsigma$, where
$s^{\rm lin}_{ij}=z_{ij}-\hat p_iy_j-\hat p_jy_i$. The scalar integrals turn
$\norm a^2+\norm{b+\varsigma c}^2$ into the quadratic form of
$\tfrac12I\oplus\bigl(\begin{smallmatrix}1/2&1/6\\1/6&1/12\end{smallmatrix}
\bigr)\otimes I$ in $(a,b,c)$; the $2\times2$ matrix has determinant $1/72$
and trace $7/12$, hence least eigenvalue at least $1/42$. With $A\succeq
\varrho_AI$ and $\norm{C_{\mathrm{dup}}s^{\rm lin}}\ge\norm{s^{\rm lin}}$, the
quadratic part is at least $\frac{\varrho_A}{42}(\norm y^2+\norm{s^{\rm lin}}^2)$.
The linear map $(y,z)\mapsto(y,s^{\rm lin})$ and its inverse have norm at most
$1+2\norm{\hat p}\le J_A$. Hence the quadratic-part matrix $H_{\hat p}$ satisfies
\[
 \mu_0I\preceq H_{\hat p}\preceq\Lambda_0I,\qquad
 \mu_0=\frac{\varrho_A}{48J_A^2},\qquad \Lambda_0=10\,U_A(P_A+2)^2J_A^2,
\]
the upper bound following from $\norm A\le U_A$,
$\norm{\hat p\otimes y}\le(P_A+1)\norm y$, $\norm{C_{\mathrm{dup}}}^2\le2$, and
$\norm{b+\varsigma c}^2\le2\norm b^2+2\norm c^2$. The point $w_*$ has norm at
most $K_A=2P_A^2$.

\emph{The gradient at $w_*$.} The vector
$\Phi_A(\hat p,y)=\nabla_{(y,z)}g_{\hat p}(y,(y_iy_j))$ is a polynomial in
$(\hat p,y)$ of degree bounded by an absolute constant, with rational
coefficients computed from $f$ and $A$. Let $\deg\Phi_A$ be its degree, $S_A$
the sum of the absolute values of its coefficients, and
$D_A=1+\deg\Phi_A\,S_A(P_A+2)^{\deg\Phi_A}$. At $\hat p=p$, the definition is a
homogeneous quadratic in $(d,s)$, because $f(p)=0$ and $\nabla f(p)=0$, and
$d$ and $s$ vanish at $w_*$; so $\Phi_A(p,p)=0$. On the box where all
coordinates of $\hat p$ have absolute value at most $P_A+1$, the Jacobian of
$\Phi_A$ in $\hat p$ has norm at most $D_A$, so
$\norm{\nabla g_{\hat p}(w_*)}=\norm{\Phi_A(\hat p,p)}\le D_A\norm{\hat p-p}$.

\emph{Residuals.} Let $r^{\rm gr}_{ij}=y_iy_j-z_{ij}$ ($i\le j$), and let
$\mathsf a(y,z)$ be obtained from $\nabla f(y)$ by replacing, in every cubic
monomial $y_iy_jy_k$ with $i\le j\le k$, the factor $y_jy_k$ by $z_{jk}$. Then
$\mathsf a(y,(y_iy_j))=\nabla f(y)$, and the $N$ quadratics
$\mathsf R=(\mathsf a,r^{\rm gr})$ have the single real common zero $w_*$,
because $r^{\rm gr}=0$ restricts to the graph and $\nabla f$ has the single
zero $p$. In the coordinates $(y,r^{\rm gr})$, related to $(y,z)$ by a map of
Jacobian determinant $\pm1$, one has $\mathsf a=\nabla f(y)-E(y)r^{\rm gr}$
with $E$ linear, so the Jacobian of $\mathsf R$ at $w_*$ is block triangular
with diagonal blocks $\nabla^2f(p)$ and $I$, and
$\abs{\det D\mathsf R(w_*)}=\det\nabla^2f(p)\ge\varrho_A^n$. Let
$W_0=1+\sum_j\norm{\mathsf R_j}_1$ (sums of absolute values of coefficients)
and $r_j=\mathsf R_j/(2W_0)$. The quadratic-part matrices of the $r_j$ have
norm at most $1/2$, their gradients at $w_*$ have norm at most
$\beta_A=1+K_A$, the Jacobian of $(r_j)$ at $w_*$ has norm at most $N\beta_A$,
and its determinant is at least $\varrho_A^n/(2W_0)^N$ in absolute value.
Hence $\sum_jb_jb_j^{\mathsf T}\succeq\nu_A^2I$ with
$\nu_A=\min\{1,\varrho_A^n(2W_0)^{-N}(N\beta_A)^{-(N-1)}\}$.

\emph{Assembly.} Choose a dyadic $t$ with
$\varepsilon=t^2\le\min\{1,\mu_0^2/(2N),\nu_A^2\mu_0^2/(36N(\Lambda_0+N\beta_A)^2)\}$,
compute $\hat p$ with $\norm{\hat p-p}\le\min\{1,\varepsilon/D_A\}$, and put
$g=g_{\hat p}$. Then $g$ and $r_1,\ldots,r_N$ satisfy the hypotheses of
Lemma~\ref{lem:quartic-realization} at $w_*$ in dimension $N$ with $N$
residuals, $\mu=\mu_0$, $\Lambda=\Lambda_0$, $\beta=\beta_A$ and $\nu=\nu_A$.
The lemma gives the stated properties of $B$; in particular the canonical
Hessian Gram of $B$ is computed directly from the rational factors, without a
second approximation of $p$. Every quantity above has polynomial bit length,
so the algorithm runs in polynomial time. By Theorem~\ref{thm:sos-length}, $B$
is not a sum of fewer than $N+1$ real squares.
\end{proof}
